% generated by GAPDoc2LaTeX from XML source (Frank Luebeck)
\documentclass[a4paper,11pt]{report}

\usepackage[top=37mm,bottom=37mm,left=27mm,right=27mm]{geometry}
\sloppy
\pagestyle{myheadings}
\usepackage{amssymb}
\usepackage[utf8]{inputenc}
\usepackage{makeidx}
\makeindex
\usepackage{color}
\definecolor{FireBrick}{rgb}{0.5812,0.0074,0.0083}
\definecolor{RoyalBlue}{rgb}{0.0236,0.0894,0.6179}
\definecolor{RoyalGreen}{rgb}{0.0236,0.6179,0.0894}
\definecolor{RoyalRed}{rgb}{0.6179,0.0236,0.0894}
\definecolor{LightBlue}{rgb}{0.8544,0.9511,1.0000}
\definecolor{Black}{rgb}{0.0,0.0,0.0}

\definecolor{linkColor}{rgb}{0.0,0.0,0.554}
\definecolor{citeColor}{rgb}{0.0,0.0,0.554}
\definecolor{fileColor}{rgb}{0.0,0.0,0.554}
\definecolor{urlColor}{rgb}{0.0,0.0,0.554}
\definecolor{promptColor}{rgb}{0.0,0.0,0.589}
\definecolor{brkpromptColor}{rgb}{0.589,0.0,0.0}
\definecolor{gapinputColor}{rgb}{0.589,0.0,0.0}
\definecolor{gapoutputColor}{rgb}{0.0,0.0,0.0}

%%  for a long time these were red and blue by default,
%%  now black, but keep variables to overwrite
\definecolor{FuncColor}{rgb}{0.0,0.0,0.0}
%% strange name because of pdflatex bug:
\definecolor{Chapter }{rgb}{0.0,0.0,0.0}
\definecolor{DarkOlive}{rgb}{0.1047,0.2412,0.0064}


\usepackage{fancyvrb}

\usepackage{mathptmx,helvet}
\usepackage[T1]{fontenc}
\usepackage{textcomp}


\usepackage[
            pdftex=true,
            bookmarks=true,        
            a4paper=true,
            pdftitle={Written with GAPDoc},
            pdfcreator={LaTeX with hyperref package / GAPDoc},
            colorlinks=true,
            backref=page,
            breaklinks=true,
            linkcolor=linkColor,
            citecolor=citeColor,
            filecolor=fileColor,
            urlcolor=urlColor,
            pdfpagemode={UseNone}, 
           ]{hyperref}

\newcommand{\maintitlesize}{\fontsize{50}{55}\selectfont}

% write page numbers to a .pnr log file for online help
\newwrite\pagenrlog
\immediate\openout\pagenrlog =\jobname.pnr
\immediate\write\pagenrlog{PAGENRS := [}
\newcommand{\logpage}[1]{\protect\write\pagenrlog{#1, \thepage,}}
%% were never documented, give conflicts with some additional packages

\newcommand{\GAP}{\textsf{GAP}}

%% nicer description environments, allows long labels
\usepackage{enumitem}
\setdescription{style=nextline}

%% depth of toc
\setcounter{tocdepth}{1}





%% command for ColorPrompt style examples
\newcommand{\gapprompt}[1]{\color{promptColor}{\bfseries #1}}
\newcommand{\gapbrkprompt}[1]{\color{brkpromptColor}{\bfseries #1}}
\newcommand{\gapinput}[1]{\color{gapinputColor}{#1}}


\begin{document}

\logpage{[ 0, 0, 0 ]}
\begin{titlepage}
\mbox{}\vfill

\begin{center}{\maintitlesize \textbf{ FusionSystems \\
\mbox{}}}\\
\vfill

\hypersetup{pdftitle={ FusionSystems }}
\markright{\scriptsize \mbox{}\hfill  FusionSystems  \hfill\mbox{}}
{\Huge \textbf{ Fusion systems in \textsf{GAP} \\
\mbox{}}}\\
\vfill

{\Huge  0.1 \mbox{}}\\[1cm]
{ 21 October 2026 \mbox{}}\\[1cm]
\mbox{}\\[2cm]
{\Large \textbf{\strut  Pete Gautam     \strut\mbox{}}}\\
{\Large \textbf{\strut  Justin Lynd     \strut\mbox{}}}\\
\hypersetup{pdfauthor={ Pete Gautam    ;  Justin Lynd    }}
\end{center}\vfill

\mbox{}\\
{\mbox{}\\
\small \noindent \textbf{ Pete Gautam    }  Email: \href{mailto://pratyush.gautam@manchester.ac.uk} {\texttt{pratyush.gautam@manchester.ac.uk}}\\
  Homepage: \href{https://github.com/pete-g00} {\texttt{https://github.com/pete\texttt{\symbol{45}}g00}}\\
  Address: \begin{minipage}[t]{8cm}\noindent
 Alan Turing Building,\\
 University of Manchester,\\
 Manchester,\\
 M13 9PY,\\
 United Kingdom\\
 \end{minipage}
}\\
{\mbox{}\\
\small \noindent \textbf{ Justin Lynd    }  Email: \href{mailto://lynd@louisiana.edu} {\texttt{lynd@louisiana.edu}}\\
  Homepage: \href{https://userweb.ucs.louisiana.edu/~jxl7583/} {\texttt{https://userweb.ucs.louisiana.edu/\texttt{\symbol{126}}jxl7583/}}\\
  Address: \begin{minipage}[t]{8cm}\noindent
 Mathematics Department\\
 University of Louisiana at Lafayette\\
 217 Maxim Doucet Hall, P.O. Box 43568\\
 Lafayette, LA 70504\texttt{\symbol{45}}3568 USA\\
 \end{minipage}
}\\
\end{titlepage}

\newpage\setcounter{page}{2}
{\small 
\section*{Abstract}
\logpage{[ 0, 0, 1 ]}
 The FusionSystems package allows for computation with fusion systems.
Currently, it supports finding proto\texttt{\symbol{45}}essential subgroups
for a given p\texttt{\symbol{45}}group and checking whether it could support
reduced fusion systems. \mbox{}}\\[1cm]
{\small 
\section*{Copyright}
\logpage{[ 0, 0, 2 ]}
 FusionSystems is free software; you can redistribute and/or modify it under
the terms of the \href{ https://www.fsf.org/licenses/gpl.html} {GNU General Public License} as published by the Free Software Foundation; either version 2 of the License,
or (at your option) any later version. \mbox{}}\\[1cm]
{\small 
\section*{Acknowledgements}
\logpage{[ 0, 0, 3 ]}
 We thank Martin van Beek for helping expand the list of
proto\texttt{\symbol{45}}essential and reduced checks. We also thank David
Burrell and Josh Bridges for their contributions to the project. \mbox{}}\\[1cm]
\newpage

\def\contentsname{Contents\logpage{[ 0, 0, 4 ]}}

\tableofcontents
\newpage

     
\chapter{\textcolor{Chapter }{Introduction}}\label{Chapter_Introduction}
\logpage{[ 1, 0, 0 ]}
\hyperdef{L}{X7DFB63A97E67C0A1}{}
{
  

 The \textsf{FusionSystems} package tests whether saturated fusion systems with given properties could
exist on a finite $p$\texttt{\symbol{45}}group $S$, without constructing them. Its main functions are: 

 
\begin{itemize}
\item  \texttt{AllProtoEssentials} (\ref{AllProtoEssentials:for IsPGroup}) finds the subgroups of $S$ that could be essential in a saturated fusion system on $S$, up to $S$\texttt{\symbol{45}}conjugacy, and \texttt{IsProtoEssential} (\ref{IsProtoEssential:for IsPGroup, IsPGroup}) tests a single subgroup (Chapters \ref{Chapter_ProtoEssentialSubgroups} and \ref{Chapter_ProtoEssentialChecks}); 
\item  \texttt{PossiblySupportsReducedFusionSystems} (\ref{PossiblySupportsReducedFusionSystems:for IsPGroup}) tests whether $S$ could support a reduced fusion system (Chapter \ref{Chapter_Checks_for_reduced_fusion_systems}); 
\item  \texttt{FindReducedSequences} (\ref{FindReducedSequences:for IsPGroup, IsList}) finds the automizer sequences that could lead to reduced fusion systems
(Chapter \ref{Chapter_AutomizerSequences}). 
\end{itemize}
 

 The tests check necessary conditions only: a subgroup found to be
proto\texttt{\symbol{45}}essential, or a group that passes the tests for
reduced fusion systems, need not lead to an actual saturated fusion system. 

 To follow the progress of a longer computation, set the info level of \texttt{InfoFusion} to 1 or higher. 

 
\begin{Verbatim}[commandchars=!@|,fontsize=\small,frame=single,label=Example]
  !gapprompt@gap>| !gapinput@S := SmallGroup(2^7, 931);;|
  !gapprompt@gap>| !gapinput@Length(AllProtoEssentials(S));|
  5
  !gapprompt@gap>| !gapinput@PossiblySupportsReducedFusionSystems(S);|
  true
\end{Verbatim}
 }

   
\chapter{\textcolor{Chapter }{Proto\texttt{\symbol{45}}essential subgroups}}\label{Chapter_ProtoEssentialSubgroups}
\logpage{[ 2, 0, 0 ]}
\hyperdef{L}{X87BC4E4D835004DD}{}
{
  
\section{\textcolor{Chapter }{Brief description of the algorithm}}\label{Chapter_Proto-essential_subgroups_Section_Brief_description_of_the_algorithm}
\logpage{[ 2, 1, 0 ]}
\hyperdef{L}{X86D9F65382E5F487}{}
{
  Let $S$ be a finite $p$\texttt{\symbol{45}}group. Recall that $E$ is proto\texttt{\symbol{45}}essential in $S$ if there \emph{could} exist a saturated fusion system $\mathcal{F}$ on $S$ such that $E \in \mathcal{E}(\mathcal{F})$. See Chapter \ref{Chapter_ProtoEssentialChecks} for the list of specific tests currently implemented. In this section, we
describe the algorithm for constructing all the
proto\texttt{\symbol{45}}essential subgroups of $S$. This is based on \cite[Section 3]{sporadics2}. We fix a central series of $S$: 
\[1 = Z_0 \leq Z_1 \leq \dots \leq Z_n = S,\]
 where $|Z_i| = p^i$ for $1 \leq i \leq n$, and there exists some $1 \leq j \leq n$ such that $Z_j = S'$. For $0 \leq i < j$, let $\pi_i \colon S \to S/Z_i$ be the projection map, and 
\[\mathfrak{C}_i := \{E \mid E = \pi^{-1}(C_S(x)), x \in S/Z_i \textrm{ of order
} p \textrm{ with } E \textrm{ proto-essential in } S\}.\]
 We let $\mathfrak{C} := \bigcup_{i=0}^j \mathfrak{C}_i$. In \cite[Theorem 3.6]{sporadics2}, Gautam shows that if $\mathcal{F}$ is a saturated fusion system on $S$ with $E \in \mathcal{E}(\mathcal{F})$, then $E$ is $\mathcal{F}$\texttt{\symbol{45}}conjugate to some $X \in \mathfrak{C}$. }

 
\section{\textcolor{Chapter }{The operations}}\label{Chapter_Proto-essential_subgroups_Section_The_operations}
\logpage{[ 2, 2, 0 ]}
\hyperdef{L}{X7CB2CAF3805CB749}{}
{
  

\subsection{\textcolor{Chapter }{AllProtoEssentials (for IsPGroup)}}
\logpage{[ 2, 2, 1 ]}\nobreak
\hyperdef{L}{X7CD47D387D876A60}{}
{\noindent\textcolor{FuncColor}{$\triangleright$\enspace\texttt{AllProtoEssentials({\mdseries\slshape S})\index{AllProtoEssentials@\texttt{AllProtoEssentials}!for IsPGroup}
\label{AllProtoEssentials:for IsPGroup}
}\hfill{\scriptsize (attribute)}}\\
\textbf{\indent Returns:\ }
a list 



 The attribute \texttt{AllProtoEssentials} constructs all proto\texttt{\symbol{45}}essential subgroups of a $p$\texttt{\symbol{45}}group $S$. It works as follows: 
\begin{itemize}
\item  it calls \texttt{MainProtoEssentials} (\ref{MainProtoEssentials}) which constructs the set $\mathfrak{C}$; 
\item  it then calls \texttt{GenerateProtoEssentials} (\ref{GenerateProtoEssentials:for IsPGroup, IsList}), which uses $\mathfrak{C}$ to construct all proto\texttt{\symbol{45}}essential subgroups, switching from $\mathcal{F}$\texttt{\symbol{45}}conjugacy to $S$\texttt{\symbol{45}}conjugacy (independent of $\mathcal{F}$). 
\end{itemize}
 

 It should be pointed out that there is no known cases where \texttt{GenerateProtoEssentials} generates a further ($S$\texttt{\symbol{45}}conjugacy) class of proto\texttt{\symbol{45}}essential
subgroups. Finds all the proto\texttt{\symbol{45}}essential subgroups of $S$, up to $S$\texttt{\symbol{45}}conjugacy. }

 
\begin{Verbatim}[commandchars=!@|,fontsize=\small,frame=single,label=Example]
  !gapprompt@gap>| !gapinput@S := SmallGroup(3^4, 7);|
  <pc group of size 81 with 4 generators>
  !gapprompt@gap>| !gapinput@L := AllProtoEssentials(S);|
  [ Group([ f4, f1, f3 ]), Group([ f4, f2, f3 ]), Group([ f1, f4 ]) ]
  !gapprompt@gap>| !gapinput@List(L, StructureDescription);|
  [ "(C3 x C3) : C3", "C3 x C3 x C3", "C3 x C3" ]
\end{Verbatim}
 

\subsection{\textcolor{Chapter }{MainProtoEssentials}}
\logpage{[ 2, 2, 2 ]}\nobreak
\hyperdef{L}{X86B32ED6791EF004}{}
{\noindent\textcolor{FuncColor}{$\triangleright$\enspace\texttt{MainProtoEssentials({\mdseries\slshape S[, i]})\index{MainProtoEssentials@\texttt{MainProtoEssentials}}
\label{MainProtoEssentials}
}\hfill{\scriptsize (function)}}\\
\textbf{\indent Returns:\ }
a list 



 Returns the list of the \texttt{main} proto\texttt{\symbol{45}}essential subgroups of $S$. A proto\texttt{\symbol{45}}essential subgroup is called \emph{main} if it lies in $\mathfrak{C}$. 

 If the second argument \texttt{i} is provided, the function behaves as follows: 
\begin{itemize}
\item  $i = 0$ \texttt{\symbol{45}}\texttt{\symbol{45}} we only compute $\mathfrak{C}_0$; 
\item  $i = -1$ \texttt{\symbol{45}}\texttt{\symbol{45}} we only compute $\mathfrak{C} \setminus \mathfrak{C}_0$; and 
\item  $i = 1$ \texttt{\symbol{45}}\texttt{\symbol{45}} we compute all $\mathfrak{C}$. 
\end{itemize}
 

 If the argument is not provided, we compute all $\mathfrak{C}$. Running \texttt{GenerateProtoEssentials}on the list $L$ returned generates all the proto\texttt{\symbol{45}}essentials subgroups of $S$. }

 
\begin{Verbatim}[commandchars=!@|,fontsize=\small,frame=single,label=Example]
  !gapprompt@gap>| !gapinput@S := SmallGroup(2^7, 931);|
  <pc group of size 128 with 7 generators>
  !gapprompt@gap>| !gapinput@C0 := MainProtoEssentials(S, 0);|
  [ Group([ f1, f2, f4, f6, f7 ]), Group([ f1, f3, f5, f6, f7 ]), 
  Group([ f2, f3, f4, f5, f6, f7 ]) ]
  !gapprompt@gap>| !gapinput@C := MainProtoEssentials(S, -1);|
  [ Group([ f7, f1, f2, f4, f5, f6 ]), Group([ f7, f1, f3, f4, f5, f6 ]) ]
\end{Verbatim}
 

\subsection{\textcolor{Chapter }{GenerateProtoEssentials (for IsPGroup, IsList)}}
\logpage{[ 2, 2, 3 ]}\nobreak
\hyperdef{L}{X86B251297A8CF204}{}
{\noindent\textcolor{FuncColor}{$\triangleright$\enspace\texttt{GenerateProtoEssentials({\mdseries\slshape S, L})\index{GenerateProtoEssentials@\texttt{GenerateProtoEssentials}!for IsPGroup, IsList}
\label{GenerateProtoEssentials:for IsPGroup, IsList}
}\hfill{\scriptsize (operation)}}\\
\textbf{\indent Returns:\ }
a list 



 Uses the list $L$ containing the main proto\texttt{\symbol{45}}essential subgroups of $S$ (i.e. the set $\mathfrak{C}$) to generate all the proto\texttt{\symbol{45}}essential subgroups of $S$. }

 
\begin{Verbatim}[commandchars=!@|,fontsize=\small,frame=single,label=Example]
  !gapprompt@gap>| !gapinput@S := SmallGroup(3^6, 321);|
  <pc group of size 729 with 6 generators>
  !gapprompt@gap>| !gapinput@L := MainProtoEssentials(S);|
  [ Group([ f6, f2, f3, f4, f5 ]), Group([ f6, f1*f2^2, f4, f5 ]) ]
  !gapprompt@gap>| !gapinput@L := GenerateProtoEssentials(S, L);|
  [ Group([ f6, f2, f3, f4, f5 ]), Group([ f6, f1*f2^2, f4, f5 ]) ]
\end{Verbatim}
 }

 }

   
\chapter{\textcolor{Chapter }{Checks for proto\texttt{\symbol{45}}essentials}}\label{Chapter_ProtoEssentialChecks}
\logpage{[ 3, 0, 0 ]}
\hyperdef{L}{X7A73DFD382D9C7CF}{}
{
  In this chapter, we describe all the tests currently implemented in GAP that
test whether a subgroup $E$ of a $p$\texttt{\symbol{45}}group $S$ is proto\texttt{\symbol{45}}essential. Recall that $E$ is proto\texttt{\symbol{45}}essential in $S$ if there \emph{could} exist a saturated fusion system $\mathcal{F}$ on $S$ such that $E \in \mathcal{E}(\mathcal{F})$. 
\section{\textcolor{Chapter }{The main functions}}\label{Chapter_Checks_for_proto-essentials_Section_The_main_functions}
\logpage{[ 3, 1, 0 ]}
\hyperdef{L}{X7D3DC4ED855DC13C}{}
{
  

\subsection{\textcolor{Chapter }{IsProtoEssential (for IsPGroup, IsPGroup)}}
\logpage{[ 3, 1, 1 ]}\nobreak
\hyperdef{L}{X7957360F807FC82A}{}
{\noindent\textcolor{FuncColor}{$\triangleright$\enspace\texttt{IsProtoEssential({\mdseries\slshape S, E})\index{IsProtoEssential@\texttt{IsProtoEssential}!for IsPGroup, IsPGroup}
\label{IsProtoEssential:for IsPGroup, IsPGroup}
}\hfill{\scriptsize (operation)}}\\
\textbf{\indent Returns:\ }
true or false 



 

 Let $S$ be a $p$\texttt{\symbol{45}}group and $E$ be a subgroup of $S$. The operation \texttt{IsProtoEssential} checks whether $E$ is a proto\texttt{\symbol{45}}essential subgroup of $S$. 

 This function checks whether the group $E$ satisfies certain properties with respect to $S$ that are known to hold when $E$ is an essential subgroup in some fusion system $\mathcal{F}$ on $E$. In this process, we construct (in order): 
\begin{itemize}
\item  $C_S(E)$; 
\item  $N_S(E)$; 
\item  $\mathrm{Aut}(N_S(E))$; and 
\item  $\mathrm{Aut}(E)$. 
\end{itemize}
 

 Each computation is considered more expensive than the previous one. See
Section \ref{Section_PEChecks} for the list of tests currently implemented. 

 Note that it is possible that $E$ is found to be proto\texttt{\symbol{45}}essential even though there is no
fusion system $\mathcal{F}$ on $S$ with $E \in \mathcal{E}(\mathcal{F})$. This can be for various reasons, for example, $E$ could be essential in $\mathcal{F}$ only if there is an overgroup of $E$ that is also essential, but there is no such group. Some of these cases may be
ruled out using automizer sequences, but it cannot be fully determined without
constructing the actual fusion systems. 

 }

 
\begin{Verbatim}[commandchars=!@|,fontsize=\small,frame=single,label=Example]
  !gapprompt@gap>| !gapinput@S := SmallGroup(3^4, 7);|
  <pc group of size 81 with 4 generators>
  !gapprompt@gap>| !gapinput@C := SubgroupsSolvableGroup(S);;|
  !gapprompt@gap>| !gapinput@C := Filtered(C, A -> IsProtoEssential(S, A));|
  [ Group([ f3, f4, f1 ]), Group([ f1, f4 ]), Group([ f3, f4, f2 ]) ]
  !gapprompt@gap>| !gapinput@List(C, StructureDescription);|
  [ "(C3 x C3) : C3", "C3 x C3", "C3 x C3 x C3" ]
\end{Verbatim}
 

\subsection{\textcolor{Chapter }{IsProtoEssentialNC (for IsPGroup, IsPGroup)}}
\logpage{[ 3, 1, 2 ]}\nobreak
\hyperdef{L}{X7808730C7E611205}{}
{\noindent\textcolor{FuncColor}{$\triangleright$\enspace\texttt{IsProtoEssentialNC({\mdseries\slshape S, E})\index{IsProtoEssentialNC@\texttt{IsProtoEssentialNC}!for IsPGroup, IsPGroup}
\label{IsProtoEssentialNC:for IsPGroup, IsPGroup}
}\hfill{\scriptsize (operation)}}\\
\textbf{\indent Returns:\ }
true or false 



 

 The operation \texttt{IsProtoEssentialNC} behaves like \texttt{IsProtoEssential}, but further assumes that: 
\begin{itemize}
\item  $E$ is a proper subgroup of $S$; 
\item  $C_S(E) \leq E$; 
\item  $E$ passes \texttt{Zed0Test}. 
\end{itemize}
 

 }

 }

 
\section{\textcolor{Chapter }{The specific functions}}\label{Section_PEChecks}
\logpage{[ 3, 2, 0 ]}
\hyperdef{L}{X85683E4783BE234C}{}
{
  

\subsection{\textcolor{Chapter }{Zed0 (for IsPGroup)}}
\logpage{[ 3, 2, 1 ]}\nobreak
\hyperdef{L}{X7DBEBC14800167CD}{}
{\noindent\textcolor{FuncColor}{$\triangleright$\enspace\texttt{Zed0({\mdseries\slshape S})\index{Zed0@\texttt{Zed0}!for IsPGroup}
\label{Zed0:for IsPGroup}
}\hfill{\scriptsize (attribute)}}\\
\textbf{\indent Returns:\ }
a subgroup 



 

 Let $S$ be a $p$\texttt{\symbol{45}}group. The attribute \texttt{Zed0} computes the subgroup $Z_0$ of $S$, which is given by: 
\[Z_0 := \langle x \in Z_2(S) \mid |[S, x]| \leq p \rangle.\]
 

 }

 

\subsection{\textcolor{Chapter }{Zed0Test (for IsPGroup, IsPGroup)}}
\logpage{[ 3, 2, 2 ]}\nobreak
\hyperdef{L}{X820687857C6BCC41}{}
{\noindent\textcolor{FuncColor}{$\triangleright$\enspace\texttt{Zed0Test({\mdseries\slshape S, E})\index{Zed0Test@\texttt{Zed0Test}!for IsPGroup, IsPGroup}
\label{Zed0Test:for IsPGroup, IsPGroup}
}\hfill{\scriptsize (operation)}}\\
\textbf{\indent Returns:\ }
true or false 



 

 Let $S$ be a $p$\texttt{\symbol{45}}group and $E$ be a subgroup of $S$. The operation \texttt{Zed0Test} checks that one of the following holds: 
\begin{itemize}
\item  $Z_0 \leq E$; or 
\item  $[N_S(E) : E] = p$, and there exists some $x \in \Omega(Z(E))$ such that $E = C_S(x)$. 
\end{itemize}
 

 Note that this function is not used when constructing $\mathfrak{C}$ as the choice of $x$ is easier to find while the set $\mathfrak{C}_0$ is being constructed. 

 This test applies at any prime $p$. It is based on \cite[Proposition 2.3(a)]{AOV17} (proven for $p=2$). 

 }

 
\begin{Verbatim}[commandchars=!@|,fontsize=\small,frame=single,label=Example]
  !gapprompt@gap>| !gapinput@S := SmallGroup(2^7, 931);|
  <pc group of size 128 with 7 generators>
  !gapprompt@gap>| !gapinput@Z0 := Zed0(S);|
  Group([ f6, f7 ])
  !gapprompt@gap>| !gapinput@C := SubgroupsSolvableGroup(S);;|
  !gapprompt@gap>| !gapinput@Length(C);|
  135
  !gapprompt@gap>| !gapinput@Number(C, A -> Zed0Test(S,A));|
  66
\end{Verbatim}
 

\subsection{\textcolor{Chapter }{RankTest (for IsPGroup, IsPGroup)}}
\logpage{[ 3, 2, 3 ]}\nobreak
\hyperdef{L}{X7CC75ED57BD90845}{}
{\noindent\textcolor{FuncColor}{$\triangleright$\enspace\texttt{RankTest({\mdseries\slshape S, E})\index{RankTest@\texttt{RankTest}!for IsPGroup, IsPGroup}
\label{RankTest:for IsPGroup, IsPGroup}
}\hfill{\scriptsize (operation)}}\\
\textbf{\indent Returns:\ }
a number 



 

 Let $S$ be a $p$\texttt{\symbol{45}}group and $E$ be a subgroup of $S$. The operation \texttt{RankTest} checks that $\textrm{Out}_S(E)$ is isomorphic to a Sylow $p$\texttt{\symbol{45}}subgroup of a group containing a strongly $p$\texttt{\symbol{45}}embedded subgroup. If so, we return a value specifying its
type, namely: 
\begin{itemize}
\item  $0$ \texttt{\symbol{45}}\texttt{\symbol{45}} $\textrm{Out}_S(E)$ is cyclic or quaternion; 
\item  $1$ \texttt{\symbol{45}}\texttt{\symbol{45}} $\textrm{Out}_S(E)$ is elementary abelian of order $\geq p^2$; 
\item  $2$ \texttt{\symbol{45}}\texttt{\symbol{45}} Sylow $p$\texttt{\symbol{45}}subgroup of $\textrm{PSU}_3(p^n)$; 
\item  $3$ \texttt{\symbol{45}}\texttt{\symbol{45}} Sylow $2$\texttt{\symbol{45}}subgroup of $\mathrm{Sz}(2^{2n+1})$; 
\item  $4$ \texttt{\symbol{45}}\texttt{\symbol{45}} Sylow $3$\texttt{\symbol{45}}subgroup of ${^2\mathrm{G}}_2(3^{2n+1})$; and 
\item  $5$ \texttt{\symbol{45}}\texttt{\symbol{45}} $3^{1+2}_-$ and $5^{1+2}_-$. 
\end{itemize}
 We return $-1$ if none of the cases above apply. During the process, we also check that the
subgroup $E$ has sufficiently high rank to ensure that $\mathrm{Aut}(E)$ has a possible valid section. 

 This test applies at any prime $p$. It is based on \cite[Proposition 6.9]{SAM14}. 

 }

 
\begin{Verbatim}[commandchars=!@|,fontsize=\small,frame=single,label=Example]
  !gapprompt@gap>| !gapinput@S := SmallGroup(3^6, 479);|
  <pc group of size 729 with 6 generators>
  !gapprompt@gap>| !gapinput@A := First(MaximalSubgroups(S), IsElementaryAbelian);|
  Group([ f2, f3, f4, f5, f6 ])
  !gapprompt@gap>| !gapinput@B := First(MaximalSubgroups(S), A -> Exponent(A) = 3 |
  !gapprompt@>| !gapinput@              and not IsElementaryAbelian(A));|
  Group([ f1, f3, f4, f5, f6 ])
  !gapprompt@gap>| !gapinput@C := First(MaximalSubgroups(S), A -> Exponent(A) = 9);|
  Group([ f1, f2, f3, f5, f6 ])
  !gapprompt@gap>| !gapinput@T := Intersection(A, B);|
  Group([ f3^2, f4^2, f5^2, f6^2 ])
  !gapprompt@gap>| !gapinput@U := Intersection(A, B, C);|
  Group([ f3^2, f5^2, f6^2 ])
  !gapprompt@gap>| !gapinput@RankTest(S, A);|
  0
  !gapprompt@gap>| !gapinput@RankTest(S, B);|
  0
  !gapprompt@gap>| !gapinput@RankTest(S, C);|
  0
  !gapprompt@gap>| !gapinput@RankTest(S, T);|
  1
  !gapprompt@gap>| !gapinput@RankTest(S, U);|
  -1
\end{Verbatim}
 

\subsection{\textcolor{Chapter }{ModuleTest (for IsPGroup, IsPGroup)}}
\logpage{[ 3, 2, 4 ]}\nobreak
\hyperdef{L}{X7CC5E68381F09D12}{}
{\noindent\textcolor{FuncColor}{$\triangleright$\enspace\texttt{ModuleTest({\mdseries\slshape S, E})\index{ModuleTest@\texttt{ModuleTest}!for IsPGroup, IsPGroup}
\label{ModuleTest:for IsPGroup, IsPGroup}
}\hfill{\scriptsize (operation)}}\\
\textbf{\indent Returns:\ }
true or false 



 

 Let $S$ be a $2$\texttt{\symbol{45}}group and $E$ be a subgroup of $S$. The operation \texttt{ModuleTest} checks that, if $\mathrm{Out}_S(E)$ has order $2^k$ (for $k \geq 2$), then for all $x \in \mathrm{Out}_S(E)$, if $x \neq 1$, then $[x, E/\Phi(E)]$ is not cyclic. 

 This test applies at $p=2$. It is based on \cite[Proposition 3.3(c)]{OV}. 

 }

 
\begin{Verbatim}[commandchars=!@|,fontsize=\small,frame=single,label=Example]
  !gapprompt@gap>| !gapinput@S := SmallGroup(2^7, 931);|
  <pc group of size 128 with 7 generators>
  !gapprompt@gap>| !gapinput@C := SubgroupsSolvableGroup(S);;|
  !gapprompt@gap>| !gapinput@Length(C);|
  135
  !gapprompt@gap>| !gapinput@Number(C, A -> Zed0Test(S,A) and ModuleTest(S,A));|
  24
\end{Verbatim}
 

\subsection{\textcolor{Chapter }{FrattiniTest (for IsPGroup, IsPGroup)}}
\logpage{[ 3, 2, 5 ]}\nobreak
\hyperdef{L}{X83A6E2E1864BAA7C}{}
{\noindent\textcolor{FuncColor}{$\triangleright$\enspace\texttt{FrattiniTest({\mdseries\slshape S, E})\index{FrattiniTest@\texttt{FrattiniTest}!for IsPGroup, IsPGroup}
\label{FrattiniTest:for IsPGroup, IsPGroup}
}\hfill{\scriptsize (operation)}}\\
\textbf{\indent Returns:\ }
true or false 



 

 Let $S$ be a $p$\texttt{\symbol{45}}group and $E$ be a subgroup of $S$. The operation \texttt{FrattiniTest} checks that 
\[C_{N_S(E)}(E/\Phi(E)) = E.\]
 

 This test applies at any prime $p$. It is based on \cite[Appendix A]{sporadics1}. 

 }

 
\begin{Verbatim}[commandchars=!@|,fontsize=\small,frame=single,label=Example]
  !gapprompt@gap>| !gapinput@S := SmallGroup(3^4, 7);|
  <pc group of size 81 with 4 generators>
  !gapprompt@gap>| !gapinput@A := First(MaximalSubgroups(S), A -> IsElementaryAbelian(A));|
  Group([ f2, f3, f4 ])
  !gapprompt@gap>| !gapinput@B := First(MaximalSubgroups(S), A -> Exponent(A) = 3 |
  !gapprompt@>| !gapinput@              and not IsElementaryAbelian(A));|
  Group([ f1, f3, f4 ])
  !gapprompt@gap>| !gapinput@C := First(MaximalSubgroups(S), A -> Exponent(A) = 9);|
  Group([ f1*f2^2, f3, f4 ])
  !gapprompt@gap>| !gapinput@T := Intersection(A, B);|
  Group([ f3^2, f4^2 ])
  !gapprompt@gap>| !gapinput@FrattiniTest(S, A);|
  true
  !gapprompt@gap>| !gapinput@FrattiniTest(S, B);|
  true
  !gapprompt@gap>| !gapinput@FrattiniTest(S, C);|
  true
  !gapprompt@gap>| !gapinput@FrattiniTest(S, T);|
  false
\end{Verbatim}
 

\subsection{\textcolor{Chapter }{LiftTest (for IsPGroup, IsPGroup, IsInt)}}
\logpage{[ 3, 2, 6 ]}\nobreak
\hyperdef{L}{X78A2281A78C9C48B}{}
{\noindent\textcolor{FuncColor}{$\triangleright$\enspace\texttt{LiftTest({\mdseries\slshape S, E, i})\index{LiftTest@\texttt{LiftTest}!for IsPGroup, IsPGroup, IsInt}
\label{LiftTest:for IsPGroup, IsPGroup, IsInt}
}\hfill{\scriptsize (operation)}}\\
\textbf{\indent Returns:\ }
true or false 



 

 Let $S$ be a $p$\texttt{\symbol{45}}group, $E$ be a subgroup of $S$ of type $i$ (see \texttt{RankTest} (\ref{RankTest:for IsPGroup, IsPGroup})). The operation \texttt{LiftTest} checks whether the maps in $N_{\mathrm{Out}_{\mathcal{F}}(E)}(\mathrm{Out}_S(E))$ are present as expected. In particular, there is some value $a$ known to divide $N_{\mathrm{Out}_{\mathcal{F}}(E)}(\mathrm{Out}_S(E))$ for each choice of $i$. We check that $a$ divides: 
\begin{itemize}
\item  $|\mathrm{Aut}(N_S(E))|$; 
\item  $|N_{\mathrm{Aut}(E)}(\mathrm{Aut}_S(E))|$; and 
\item  $|N_{\mathrm{Out}(E)}(\mathrm{Out}_S(E)) :
C_{\mathrm{Out}(E)}(\mathrm{Out}_S(E))|$. 
\end{itemize}
 

 This test applies at any prime $p$. It is based on \cite[Lemma A.1]{sporadics1} and \cite[Proposition 3.3(b)]{OV}. 

 }

 
\begin{Verbatim}[commandchars=!@|,fontsize=\small,frame=single,label=Example]
  !gapprompt@gap>| !gapinput@S := SmallGroup(3^6, 479);|
  <pc group of size 729 with 6 generators>
  !gapprompt@gap>| !gapinput@A := First(MaximalSubgroups(S), IsElementaryAbelian);|
  Group([ f2, f3, f4, f5, f6 ])
  !gapprompt@gap>| !gapinput@B := First(MaximalSubgroups(S), A -> Exponent(A) = 3 |
  !gapprompt@>| !gapinput@              and not IsElementaryAbelian(A));|
  Group([ f1, f3, f4, f5, f6 ])
  !gapprompt@gap>| !gapinput@C := First(MaximalSubgroups(S), A -> Exponent(A) = 9);|
  Group([ f1, f2, f3, f5, f6 ])
  !gapprompt@gap>| !gapinput@T := Intersection(A, B);|
  Group([ f3^2, f4^2, f5^2, f6^2 ])
  !gapprompt@gap>| !gapinput@LiftTest(S, A, 0);|
  true
  !gapprompt@gap>| !gapinput@LiftTest(S, B, 0);|
  true
  !gapprompt@gap>| !gapinput@LiftTest(S, C, 0);|
  true
  !gapprompt@gap>| !gapinput@LiftTest(S, T, 1);|
  false
\end{Verbatim}
 

\subsection{\textcolor{Chapter }{ValidAutTest (for IsPGroup, IsPGroup, IsInt)}}
\logpage{[ 3, 2, 7 ]}\nobreak
\hyperdef{L}{X86FDD62186F7F9CA}{}
{\noindent\textcolor{FuncColor}{$\triangleright$\enspace\texttt{ValidAutTest({\mdseries\slshape S, E, i})\index{ValidAutTest@\texttt{ValidAutTest}!for IsPGroup, IsPGroup, IsInt}
\label{ValidAutTest:for IsPGroup, IsPGroup, IsInt}
}\hfill{\scriptsize (operation)}}\\
\textbf{\indent Returns:\ }
true or false 



 

 Let $S$ be a $p$\texttt{\symbol{45}}group, $E$ be a subgroup of $S$ of type $i$ (see \texttt{RankTest} (\ref{RankTest:for IsPGroup, IsPGroup})). The operation \texttt{ValidAutTest} checks whether there is a valid section of $\mathrm{Aut}(E)$ that contains the expected $\mathrm{Aut}_{\mathcal{F}}(E)$. Note that this check is up to (order of) composition factors, and does not
imply that the expected section actually exists inside $\mathrm{Aut}(E)$. 

 This test applies at any prime $p$. It is based on \cite[Proposition 6.9]{SAM14} and \cite[Proposition 4.5]{PS_MAGMA}. 

 }

 
\begin{Verbatim}[commandchars=!@|,fontsize=\small,frame=single,label=Example]
  !gapprompt@gap>| !gapinput@S := SmallGroup(3^4, 7);|
  <pc group of size 81 with 4 generators>
  !gapprompt@gap>| !gapinput@A := First(MaximalSubgroups(S), A -> IsElementaryAbelian(A));|
  Group([ f2, f3, f4 ])
  !gapprompt@gap>| !gapinput@B := First(MaximalSubgroups(S), A -> Exponent(A) = 3 |
  !gapprompt@>| !gapinput@              and not IsElementaryAbelian(A));|
  Group([ f1, f3, f4 ])
  !gapprompt@gap>| !gapinput@C := First(MaximalSubgroups(S), A -> Exponent(A) = 9);|
  Group([ f1*f2^2, f3, f4 ])
  !gapprompt@gap>| !gapinput@ValidAutTest(S, A, 0);|
  true
  !gapprompt@gap>| !gapinput@ValidAutTest(S, B, 0);|
  true
  !gapprompt@gap>| !gapinput@ValidAutTest(S, C, 0);|
  false
\end{Verbatim}
 

\subsection{\textcolor{Chapter }{RadicalTest (for IsPGroup, IsPGroup)}}
\logpage{[ 3, 2, 8 ]}\nobreak
\hyperdef{L}{X7D2C79FB78684AC9}{}
{\noindent\textcolor{FuncColor}{$\triangleright$\enspace\texttt{RadicalTest({\mdseries\slshape S, E})\index{RadicalTest@\texttt{RadicalTest}!for IsPGroup, IsPGroup}
\label{RadicalTest:for IsPGroup, IsPGroup}
}\hfill{\scriptsize (operation)}}\\
\textbf{\indent Returns:\ }
true or false 



 

 Let $S$ be a $p$\texttt{\symbol{45}}group, $E$ be a proto\texttt{\symbol{45}}essential subgroup. The operation \texttt{RadicalTest}checks whether 

 
\[|O_p(\textrm{Aut}(E)) \cap \textrm{Aut}_S(E)| = [E : Z(E)].\]
 

 This test applies at any prime $p$. It is based on \cite[Proposition 4.8(2)]{PS_MAGMA}. 

 }

 
\begin{Verbatim}[commandchars=!@|,fontsize=\small,frame=single,label=Example]
  !gapprompt@gap>| !gapinput@S := SmallGroup(3^4, 7);|
  <pc group of size 81 with 4 generators>
  !gapprompt@gap>| !gapinput@A := First(MaximalSubgroups(S), A -> IsElementaryAbelian(A));|
  Group([ f2, f3, f4 ])
  !gapprompt@gap>| !gapinput@B := First(MaximalSubgroups(S), A -> |
  !gapprompt@>| !gapinput@              Exponent(A) = 3 and not IsElementaryAbelian(A));|
  Group([ f1, f3, f4 ])
  !gapprompt@gap>| !gapinput@C := First(MaximalSubgroups(S), A -> Exponent(A) = 9);|
  Group([ f1*f2^2, f3, f4 ])
  !gapprompt@gap>| !gapinput@RadicalTest(S, A);|
  true
  !gapprompt@gap>| !gapinput@RadicalTest(S, B);|
  true
  !gapprompt@gap>| !gapinput@RadicalTest(S, C);|
  false
\end{Verbatim}
 }

 }

   
\chapter{\textcolor{Chapter }{Checks for reduced fusion systems}}\label{Chapter_Checks_for_reduced_fusion_systems}
\logpage{[ 4, 0, 0 ]}
\hyperdef{L}{X87497673854B812F}{}
{
  In this chapter, we describe all the tests currently implemented in GAP that
test whether a group $S$ \emph{could} support reduced fusion systems. This is done without constructing the fusion
systems, and even possible automizers for any fusion system. As such, we can
determine whether there exists any reduced fusion system $\mathcal{F}$ on $S$ for a fixed choice of $\mathcal{E}(\mathcal{F})$. We can also determine this for any choice of $\mathcal{E}(\mathcal{F})$. 
\section{\textcolor{Chapter }{The main functions}}\label{Chapter_Checks_for_reduced_fusion_systems_Section_The_main_functions}
\logpage{[ 4, 1, 0 ]}
\hyperdef{L}{X7D3DC4ED855DC13C}{}
{
  

\subsection{\textcolor{Chapter }{PossiblySupportsReducedFusionSystems (for IsPGroup)}}
\logpage{[ 4, 1, 1 ]}\nobreak
\hyperdef{L}{X829638DD78A8FD08}{}
{\noindent\textcolor{FuncColor}{$\triangleright$\enspace\texttt{PossiblySupportsReducedFusionSystems({\mdseries\slshape S})\index{PossiblySupportsReducedFusionSystems@\texttt{Possibly}\-\texttt{Supports}\-\texttt{Reduced}\-\texttt{Fusion}\-\texttt{Systems}!for IsPGroup}
\label{PossiblySupportsReducedFusionSystems:for IsPGroup}
}\hfill{\scriptsize (attribute)}}\\
\textbf{\indent Returns:\ }
true or false 



 

 Given a $p$\texttt{\symbol{45}}group $S$, the attribute \texttt{PossiblySupportsReducedFusionSystems} checks whether $S$ \emph{could} support reduced fusion systems. 

 Note that it is not possible to check whether a group can support reduced
fusion systems without actually constructing them. This is because the
expected automizer sequences that lead to reduced fusion systems may not be
saturated, and the only way we can discover this is by checking some radical,
non\texttt{\symbol{45}}essential, subgroup in the expected saturated fusion
system. Nonetheless, $p$\texttt{\symbol{45}}groups that pass the tests in this chapter are quite
likely to support reduced fusion systems. 

 This function checks whether the group $S$ satisfies certain properties that are known to hold for a $p$\texttt{\symbol{45}}group that supports reduced fusion systems. In the
process, we compute the proto\texttt{\symbol{45}}essential subgroups of $S$, along with their automorphism groups. See Section \ref{Section_RedSpecificFn} for the list of specific tests currently implemented. We do not construct
automizer sequences. 

 }

 
\begin{Verbatim}[commandchars=!@|,fontsize=\small,frame=single,label=Example]
  !gapprompt@gap>| !gapinput@S := SmallGroup(2^7, 931);|
  <pc group of size 128 with 7 generators>
  !gapprompt@gap>| !gapinput@PossiblySupportsReducedFusionSystems(S);|
  true
\end{Verbatim}
 

\subsection{\textcolor{Chapter }{PossiblySupportsReducedFusionSystems{\textunderscore}PE (for IsPGroup, IsList)}}
\logpage{[ 4, 1, 2 ]}\nobreak
\hyperdef{L}{X87A01C45801E726C}{}
{\noindent\textcolor{FuncColor}{$\triangleright$\enspace\texttt{PossiblySupportsReducedFusionSystems{\textunderscore}PE({\mdseries\slshape S, L})\index{PossiblySupportsReducedFusionSystems{\textunderscore}PE@\texttt{Possibly}\-\texttt{Supports}\-\texttt{Reduced}\-\texttt{Fusion}\-\texttt{Systems{\textunderscore}PE}!for IsPGroup, IsList}
\label{PossiblySupportsReducedFusionSystemsuScorePE:for IsPGroup, IsList}
}\hfill{\scriptsize (operation)}}\\
\textbf{\indent Returns:\ }
true or false 



 

 Given a $p$\texttt{\symbol{45}}group $S$ and a list of proto\texttt{\symbol{45}}essential subgroups $L$, the operation \texttt{PossiblySupportsReducedFusionSystems{\textunderscore}PE} checks whether there \emph{could} exist a reduced fusion system $\mathcal{F}$ on $S$ where the set of essential subgroups $\mathcal{E}(\mathcal{F})$ is a subset of $L$. Note that we run \texttt{GenerateProtoEssentials} on $L$ (which can, in theory, increase the number of subgroups in $L$). 

 Note that it is not possible to check whether a group can support reduced
fusion systems without actually constructing them. This is because the
expected automizer sequences that lead to reduced fusion systems may not be
saturated, and the only way we can discover this is by checking some radical,
non\texttt{\symbol{45}}essential, subgroup in the expected saturated fusion
system. Nonetheless, $p$\texttt{\symbol{45}}groups that pass the tests in this chapter are quite
likely to support reduced fusion systems. 

 This function checks whether the group $S$ satisfies certain properties that are known to hold for a $p$\texttt{\symbol{45}}group that supports reduced fusion systems. We also
construct the automorphism group for subgroups in $L$. See Section \ref{Section_RedSpecificFn} for the list of specific tests currently implemented. We do not construct
automizer sequences. 

 }

 
\begin{Verbatim}[commandchars=!@|,fontsize=\small,frame=single,label=Example]
  !gapprompt@gap>| !gapinput@S := SmallGroup(2^7, 931);|
  <pc group of size 128 with 7 generators>
  !gapprompt@gap>| !gapinput@L := AllProtoEssentials(S);|
  [ Group([ f7, f1, f2, f4, f5, f6 ]), Group([ f5, f6, f7, f1, f3, f4 ]), 
  Group([ f5, f6, f7, f2, f3, f4 ]), Group([ f1, f2, f4, f6, f7 ]), 
  Group([ f1, f3, f5, f6, f7 ]) ]
  !gapprompt@gap>| !gapinput@PossiblySupportsReducedFusionSystems_PE(S, L{[1..3]});|
  true
  !gapprompt@gap>| !gapinput@PossiblySupportsReducedFusionSystems_PE(S, L{[3..5]});|
  false
\end{Verbatim}
 }

 
\section{\textcolor{Chapter }{The specific functions}}\label{Section_RedSpecificFn}
\logpage{[ 4, 2, 0 ]}
\hyperdef{L}{X85683E4783BE234C}{}
{
  

\subsection{\textcolor{Chapter }{IsDSWGroup (for IsPGroup)}}
\logpage{[ 4, 2, 1 ]}\nobreak
\hyperdef{L}{X83567EE27C9A327E}{}
{\noindent\textcolor{FuncColor}{$\triangleright$\enspace\texttt{IsDSWGroup({\mdseries\slshape S})\index{IsDSWGroup@\texttt{IsDSWGroup}!for IsPGroup}
\label{IsDSWGroup:for IsPGroup}
}\hfill{\scriptsize (attribute)}}\\
\textbf{\indent Returns:\ }
true or false 



 

 Given a $p$\texttt{\symbol{45}}group $S$, the attribute \texttt{IsDSWGroup} checks whether $p=2$, and that $S$ is isomorphic to dihedral, semi\texttt{\symbol{45}}dihedral group, or the
wreath product $C_{2^a} \wr C_2$. 

 Using the work of Oliver and his collaborators, there are many properties on a $2$\texttt{\symbol{45}}group $S$ that imply that, if $S$ supports reduced fusion systems, then $S$ is isomorphic to one of the three groups (and is called a $\mathcal{DSW}$\texttt{\symbol{45}}group) (see, for example, \cite[Proposition 5.2]{AOV13}). 

 }

 
\begin{Verbatim}[commandchars=!@|,fontsize=\small,frame=single,label=Example]
  !gapprompt@gap>| !gapinput@S := DihedralGroup(2^3);|
  <pc group of size 8 with 3 generators>
  !gapprompt@gap>| !gapinput@IsDSWGroup(S);|
  true
  !gapprompt@gap>| !gapinput@S := SmallGroup(2^4, 8);|
  <pc group of size 16 with 4 generators>
  !gapprompt@gap>| !gapinput@IsDSWGroup(S);|
  true
  !gapprompt@gap>| !gapinput@S := QuaternionGroup(2^5);|
  <pc group of size 32 with 5 generators>
  !gapprompt@gap>| !gapinput@IsDSWGroup(S);|
  false
  !gapprompt@gap>| !gapinput@S := WreathProduct(CyclicGroup(IsPermGroup, 2^3), Group((1,2)));|
  Group([ (1,2,3,4,5,6,7,8), (9,10,11,12,13,14,15,16), 
  (1,9)(2,10)(3,11)(4,12)(5,13)(6,14)(7,15)(8,16) ])
  !gapprompt@gap>| !gapinput@IdGroup(S);|
  [128, 67]
  !gapprompt@gap>| !gapinput@IsDSWGroup(S);|
  true
  !gapprompt@gap>| !gapinput@S := SmallGroup(2^7, 437);|
  <pc group of size 128 with 7 generators>
  !gapprompt@gap>| !gapinput@StructureDescription(S);|
  "(C8 x C8) : C2"
  !gapprompt@gap>| !gapinput@IsDSWGroup(S);|
  false
\end{Verbatim}
 

\subsection{\textcolor{Chapter }{ExponentTest (for IsPGroup)}}
\logpage{[ 4, 2, 2 ]}\nobreak
\hyperdef{L}{X7E63D2A97A637BAD}{}
{\noindent\textcolor{FuncColor}{$\triangleright$\enspace\texttt{ExponentTest({\mdseries\slshape S})\index{ExponentTest@\texttt{ExponentTest}!for IsPGroup}
\label{ExponentTest:for IsPGroup}
}\hfill{\scriptsize (attribute)}}\\
\textbf{\indent Returns:\ }
true or false 



 

 Let $S$ be a $p$\texttt{\symbol{45}}group. If $S$ has class $c$, the attribute \texttt{ExponentTest} checks that $Z(S)$ has exponent at most $p^t$, where 
\[t := \left\lceil \frac{c-1}{p-1} \right\rceil.\]
 

 This test applies at any prime $p$. It is based on \cite[Theorem A]{LP}. 

 }

 
\begin{Verbatim}[commandchars=!@|,fontsize=\small,frame=single,label=Example]
  !gapprompt@gap>| !gapinput@S := SmallGroup(3^4, 13);|
  <pc group of size 81 with 4 generators>
  !gapprompt@gap>| !gapinput@StructureDescription(S);|
  "C3 x (C9 : C3)"
  !gapprompt@gap>| !gapinput@ExponentTest(S);|
  true
  !gapprompt@gap>| !gapinput@S := SmallGroup(3^4, 14);|
  <pc group of size 81 with 4 generators>
  !gapprompt@gap>| !gapinput@StructureDescription(S);|
  "(C9 x C3) : C3"
  !gapprompt@gap>| !gapinput@ExponentTest(S);|
  false
\end{Verbatim}
 

\subsection{\textcolor{Chapter }{AbelianTest (for IsPGroup)}}
\logpage{[ 4, 2, 3 ]}\nobreak
\hyperdef{L}{X7F0F966F7C7E4B8D}{}
{\noindent\textcolor{FuncColor}{$\triangleright$\enspace\texttt{AbelianTest({\mdseries\slshape S})\index{AbelianTest@\texttt{AbelianTest}!for IsPGroup}
\label{AbelianTest:for IsPGroup}
}\hfill{\scriptsize (attribute)}}\\
\textbf{\indent Returns:\ }
true or false 



 

 Let $S$ be a $2$\texttt{\symbol{45}}group. The attribute \texttt{AbelianTest} first asserts whether there exists an abelian subgroup $A$ of $S$ satisfying $[N_S(A) : A] = 2$ and: 
\begin{itemize}
\item  $N_S(A) = S$; or 
\item  $N_S(A) < S$, with $\mathrm{Aut}_S(A) \nleq O_2(\mathrm{Aut})(A)$. 
\end{itemize}
 If so, the function returns true only if $S$ is a $\mathcal{DSW}$\texttt{\symbol{45}}group. 

 This test only applies at $p=2$. It is based on \cite[Proposition 5.2(a), (c)]{AOV13}. 

 }

 
\begin{Verbatim}[commandchars=!@|,fontsize=\small,frame=single,label=Example]
  !gapprompt@gap>| !gapinput@S := DihedralGroup(2^4);|
  <pc group of size 16 with 4 generators>
  !gapprompt@gap>| !gapinput@AbelianTest(S);|
  true
  !gapprompt@gap>| !gapinput@S := SmallGroup(2^5, 9);|
  <pc group of size 32 with 5 generators>
  !gapprompt@gap>| !gapinput@StructureDescription(S);|
  "(C8 x C2) : C2"
  !gapprompt@gap>| !gapinput@StructureDescription(Center(S));|
  "C2 x C2"
  !gapprompt@gap>| !gapinput@AbelianTest(S);|
  false
\end{Verbatim}
 

\subsection{\textcolor{Chapter }{DerivedCyclicTest (for IsPGroup)}}
\logpage{[ 4, 2, 4 ]}\nobreak
\hyperdef{L}{X782BA19382650E0B}{}
{\noindent\textcolor{FuncColor}{$\triangleright$\enspace\texttt{DerivedCyclicTest({\mdseries\slshape S})\index{DerivedCyclicTest@\texttt{DerivedCyclicTest}!for IsPGroup}
\label{DerivedCyclicTest:for IsPGroup}
}\hfill{\scriptsize (attribute)}}\\
\textbf{\indent Returns:\ }
true or false 



 

 Let $S$ be a $p$\texttt{\symbol{45}}group. If $S'$ is cyclic, then the attribute \texttt{DerivedCyclicTest} returns true only if: 
\begin{itemize}
\item  $S$ is a $\mathcal{DSW}$\texttt{\symbol{45}}group; or 
\item  $S$ has an abelian subgroup of index $p$ (at $p$ odd). 
\end{itemize}
 

 This test applies at any prime $p$. It is based on \cite[Proposition 5.2(b)]{AOV13} (if $p=2$). 

 }

 
\begin{Verbatim}[commandchars=!@|,fontsize=\small,frame=single,label=Example]
  !gapprompt@gap>| !gapinput@S := DihedralGroup(2^4);|
  <pc group of size 16 with 4 generators>
  !gapprompt@gap>| !gapinput@DerivedCyclicTest(S);|
  true
  !gapprompt@gap>| !gapinput@S := SmallGroup(2^5, 9);|
  <pc group of size 32 with 5 generators>
  !gapprompt@gap>| !gapinput@StructureDescription(S);|
  "(C8 x C2) : C2"
  !gapprompt@gap>| !gapinput@StructureDescription(Center(S));|
  "C2 x C2"
  !gapprompt@gap>| !gapinput@DerivedCyclicTest(S);|
  false
\end{Verbatim}
 

\subsection{\textcolor{Chapter }{ValidAutTest (for IsPGroup)}}
\logpage{[ 4, 2, 5 ]}\nobreak
\hyperdef{L}{X7E28CE0083D26C55}{}
{\noindent\textcolor{FuncColor}{$\triangleright$\enspace\texttt{ValidAutTest({\mdseries\slshape S})\index{ValidAutTest@\texttt{ValidAutTest}!for IsPGroup}
\label{ValidAutTest:for IsPGroup}
}\hfill{\scriptsize (attribute)}}\\
\textbf{\indent Returns:\ }
true or false 



 

 Let $S$ be a $p$\texttt{\symbol{45}}group (for $p$ odd). The attribute \texttt{ValidAutTest} returns true only if $\mathrm{Aut}(S)$ is not a $p$\texttt{\symbol{45}}group. If we further know that $S$ has class $\leq p-1$, and $\mathrm{Aut}(S)$ has a Hall $p$\texttt{\symbol{45}}complement $H$, we check that $[S, H] = S$. 

 This test does not apply if $p \leq 3$. It is based on \cite[7.82]{CRA} and \cite{DGPS}. 

 }

 
\begin{Verbatim}[commandchars=!@|,fontsize=\small,frame=single,label=Example]
  !gapprompt@gap>| !gapinput@S := SmallGroup(5^5, 38);|
  <pc group of size 3125 with 5 generators>
  !gapprompt@gap>| !gapinput@NilpotencyClassOfGroup(S);|
  4
  !gapprompt@gap>| !gapinput@AutomorphismGroup(S);|
  <group of size 15625 with 6 generators>
  !gapprompt@gap>| !gapinput@ValidAutTest(S);|
  false
  !gapprompt@gap>| !gapinput@S := SmallGroup(5^5, 30);|
  <pc group of size 3125 with 5 generators>
  !gapprompt@gap>| !gapinput@NilpotencyClassOfGroup(S);|
  4
  !gapprompt@gap>| !gapinput@AutomorphismGroup(S);|
  <group of size 1250000 with 11 generators>
  !gapprompt@gap>| !gapinput@ValidAutTest(S);|
  true
  !gapprompt@gap>| !gapinput@S := SmallGroup(5^5, 32);|
  <pc group of size 3125 with 5 generators>
  !gapprompt@gap>| !gapinput@NilpotencyClassOfGroup(S);|
  4
  !gapprompt@gap>| !gapinput@AutomorphismGroup(S);|
  <group of size 62500 with 8 generators>
  !gapprompt@gap>| !gapinput@ValidAutTest(S);|
  false
\end{Verbatim}
 

\subsection{\textcolor{Chapter }{OmegaAbelianizationTest (for IsPGroup)}}
\logpage{[ 4, 2, 6 ]}\nobreak
\hyperdef{L}{X852FEF1D87916D36}{}
{\noindent\textcolor{FuncColor}{$\triangleright$\enspace\texttt{OmegaAbelianizationTest({\mdseries\slshape S})\index{OmegaAbelianizationTest@\texttt{OmegaAbelianizationTest}!for IsPGroup}
\label{OmegaAbelianizationTest:for IsPGroup}
}\hfill{\scriptsize (attribute)}}\\
\textbf{\indent Returns:\ }
true or false 



 

 Let $S$ be a $2$\texttt{\symbol{45}}group, $D := S'$ and $O := \Omega(Z(S))$. The attribute \texttt{OmegaAbelianizationTest} returns false only if: 
\begin{itemize}
\item  $[\Omega(Z(S)) : \Omega(Z(S)) \cap D] = 2$; or 
\item  $\Omega(Z(S)) \nleq D$ if $\mathrm{Aut}(S)$ is a $2$\texttt{\symbol{45}}group. 
\end{itemize}
 

 This test only applies at $p=2$. It is based on \cite[Corollary I.8.5]{AKO}. 

 }

 
\begin{Verbatim}[commandchars=!@|,fontsize=\small,frame=single,label=Example]
  !gapprompt@gap>| !gapinput@S := DihedralGroup(2^4);|
  <pc group of size 16 with 4 generators>
  !gapprompt@gap>| !gapinput@OmegaAbelianizationTest(S);|
  true
  !gapprompt@gap>| !gapinput@S := SmallGroup(2^4, 3);|
  <pc group of size 16 with 4 generators>
  !gapprompt@gap>| !gapinput@StructureDescription(S);|
  "(C4 x C2) : C2"
  !gapprompt@gap>| !gapinput@OmegaAbelianizationTest(S);|
  false
\end{Verbatim}
 
\subsection{\textcolor{Chapter }{Transfer test for a group}}\label{AutoDoc_generated_group1}
\logpage{[ 4, 2, 7 ]}
\hyperdef{L}{X82E6548F81B8F05C}{}
{
\noindent\textcolor{FuncColor}{$\triangleright$\enspace\texttt{TransferMaximalTest({\mdseries\slshape S})\index{TransferMaximalTest@\texttt{TransferMaximalTest}!for IsPGroup}
\label{TransferMaximalTest:for IsPGroup}
}\hfill{\scriptsize (attribute)}}\\
\noindent\textcolor{FuncColor}{$\triangleright$\enspace\texttt{TransferFpMaximalTest({\mdseries\slshape S})\index{TransferFpMaximalTest@\texttt{TransferFpMaximalTest}!for IsPGroup}
\label{TransferFpMaximalTest:for IsPGroup}
}\hfill{\scriptsize (attribute)}}\\
\textbf{\indent Returns:\ }
true or false 



 

 Let $S$ be a $2$\texttt{\symbol{45}}group. The attribute \texttt{TransferMaximalTest} checks that, if $\mathrm{Aut}(S)$ is a $2$\texttt{\symbol{45}}group, then 
\[\bigcap_{M \in \mathcal{M}(S)} \ker(\mathrm{tr}_M^S) = S',\]
 where $\mathcal{M}(S)$ is the set of maximal subgroups of $S$, and $\mathrm{tr}_M^S \colon S \to M/M'$ is the transfer homomorphism. Alternatively, \texttt{TransferFpMaximalTest} checks that, if $\mathrm{Aut}(S)$ is a $2$\texttt{\symbol{45}}group, then 
\[\bigcap_{M \in \mathcal{M}(S)} \ker(\mathrm{tr'}_M^S) = \Phi(S),\]
 where $\mathcal{M}(S)$ is the set of maximal subgroups of $S$, and $\mathrm{tr'}_M^S \colon S \to M/\Phi(M)$ is the transfer homomorphism. 

 These tests only apply at $p=2$ when $\mathrm{Aut}(S)$ is a $2$\texttt{\symbol{45}}group. It is based on \cite[Lemma 2.1]{AOV17}. 

 }

 
\begin{Verbatim}[commandchars=!@|,fontsize=\small,frame=single,label=Example]
  !gapprompt@gap>| !gapinput@S := SmallGroup(2^7, 903);|
  <pc group of size 128 with 7 generators>
  !gapprompt@gap>| !gapinput@TransferMaximalTest(S);|
  false
  !gapprompt@gap>| !gapinput@TransferFpMaximalTest(S);|
  true
\end{Verbatim}
 
\subsection{\textcolor{Chapter }{Transfer test with a list}}\label{AutoDoc_generated_group2}
\logpage{[ 4, 2, 8 ]}
\hyperdef{L}{X7ECB69FB82C13217}{}
{
\noindent\textcolor{FuncColor}{$\triangleright$\enspace\texttt{TransferTest({\mdseries\slshape S, L})\index{TransferTest@\texttt{TransferTest}!for IsPGroup, IsList}
\label{TransferTest:for IsPGroup, IsList}
}\hfill{\scriptsize (operation)}}\\
\noindent\textcolor{FuncColor}{$\triangleright$\enspace\texttt{TransferFpTest({\mdseries\slshape S, L})\index{TransferFpTest@\texttt{TransferFpTest}!for IsPGroup, IsList}
\label{TransferFpTest:for IsPGroup, IsList}
}\hfill{\scriptsize (operation)}}\\
\textbf{\indent Returns:\ }
true or false 



 

 Let $S$ be a $2$\texttt{\symbol{45}}group and $L$ be a list of (proto\texttt{\symbol{45}}essential) subgroups of $S$. The operation \texttt{TransferTest} checks that, if $\mathrm{Aut}(S)$ is a $2$\texttt{\symbol{45}}group, then 
\[\bigcap_{M \in L} \ker(\mathrm{tr}_M^S) = S',\]
 where $\mathrm{tr}_M^S \colon S \to M/M'$ is the transfer homomorphism. Alternatively, \texttt{TransferFpTest} checks that, if $\mathrm{Aut}(S)$ is a $2$\texttt{\symbol{45}}group, then 
\[\bigcap_{M \in L} \ker(\mathrm{tr'}_M^S) = \Phi(S),\]
 where $\mathrm{tr'}_M^S \colon S \to M/\Phi(M)$ is the transfer homomorphism. 

 These tests only apply at $p=2$ when $\mathrm{Aut}(S)$ is a $2$\texttt{\symbol{45}}group. It is based on \cite[Lemma 2.1]{AOV17}. 

 }

 
\begin{Verbatim}[commandchars=!@|,fontsize=\small,frame=single,label=Example]
  !gapprompt@gap>| !gapinput@S := SmallGroup(2^7, 931);|
  <pc group of size 128 with 7 generators>
  !gapprompt@gap>| !gapinput@L := AllProtoEssentials(S);|
  [ Group([ f7, f1, f2, f4, f5, f6 ]), Group([ f5, f6, f7, f1, f3, f4 ]), 
  Group([ f5, f6, f7, f2, f3, f4 ]), Group([ f1, f2, f4, f6, f7 ]), 
  Group([ f1, f3, f5, f6, f7 ]) ]
  !gapprompt@gap>| !gapinput@TransferTest(S, L{[1..3]});|
  true
  !gapprompt@gap>| !gapinput@TransferTest(S, L{[3..5]});|
  false
\end{Verbatim}
 

\subsection{\textcolor{Chapter }{C0Test (for IsPGroup, IsList)}}
\logpage{[ 4, 2, 9 ]}\nobreak
\hyperdef{L}{X8660B3728652622E}{}
{\noindent\textcolor{FuncColor}{$\triangleright$\enspace\texttt{C0Test({\mdseries\slshape S, L})\index{C0Test@\texttt{C0Test}!for IsPGroup, IsList}
\label{C0Test:for IsPGroup, IsList}
}\hfill{\scriptsize (operation)}}\\
\textbf{\indent Returns:\ }
true or false 



 

 Let $S$ be a $p$\texttt{\symbol{45}}group, and $L := \mathfrak{C}_0$. The operation \texttt{C0Test} checks that there exists some subgroup $E$ in $L$ such that $\Omega(Z(S))$ is not semi\texttt{\symbol{45}}characteristic in $E$, i.e. invariant under $O^p(O^{p'}(\mathrm{Aut}(E)))$. 

 This test applies at any prime $p$. It is based on \cite[Corollary 3.12]{sporadics2}. 

 }

 
\begin{Verbatim}[commandchars=!@|,fontsize=\small,frame=single,label=Example]
  !gapprompt@gap>| !gapinput@L := AllSmallGroups(3^5, A -> NilpotencyClassOfGroup(A) >= 3);;|
  !gapprompt@gap>| !gapinput@Length(L);|
  32
  !gapprompt@gap>| !gapinput@Number(L, S -> C0Test(S, MainProtoEssentials(S, 0)));|
  17
\end{Verbatim}
 

\subsection{\textcolor{Chapter }{CTest (for IsPGroup, IsList)}}
\logpage{[ 4, 2, 10 ]}\nobreak
\hyperdef{L}{X823B722F86499207}{}
{\noindent\textcolor{FuncColor}{$\triangleright$\enspace\texttt{CTest({\mdseries\slshape S, L})\index{CTest@\texttt{CTest}!for IsPGroup, IsList}
\label{CTest:for IsPGroup, IsList}
}\hfill{\scriptsize (operation)}}\\
\textbf{\indent Returns:\ }
true or false 



 

 Let $S$ be a $p$\texttt{\symbol{45}}group, and let $L$ be a list of proto\texttt{\symbol{45}}essential subgroups of $S$. The operation \texttt{CTest} checks that: 
\begin{itemize}
\item  $L$ has at least $2$ elements, or 
\item  $L$ consists of a single subgroup, which is not normal in $S$, and $\mathrm{Aut}(S)$ is not a $p$\texttt{\symbol{45}}group. 
\end{itemize}
 

 This test applies at any prime $p$. It is based on \cite[Lemma 2.2(e)]{AOV17}. 

 }

 
\begin{Verbatim}[commandchars=!@|,fontsize=\small,frame=single,label=Example]
  !gapprompt@gap>| !gapinput@L := AllSmallGroups(3^5, A -> NilpotencyClassOfGroup(A) >= 3);;|
  !gapprompt@gap>| !gapinput@Length(L);|
  32
  !gapprompt@gap>| !gapinput@Number(L, S -> CTest(S, AllProtoEssentials(S)));|
  16
\end{Verbatim}
 

\subsection{\textcolor{Chapter }{SemiCharacteristicGeneratedTest (for IsPGroup, IsList)}}
\logpage{[ 4, 2, 11 ]}\nobreak
\hyperdef{L}{X830693B38724090D}{}
{\noindent\textcolor{FuncColor}{$\triangleright$\enspace\texttt{SemiCharacteristicGeneratedTest({\mdseries\slshape S, L})\index{SemiCharacteristicGeneratedTest@\texttt{SemiCharacteristicGeneratedTest}!for IsPGroup, IsList}
\label{SemiCharacteristicGeneratedTest:for IsPGroup, IsList}
}\hfill{\scriptsize (operation)}}\\
\textbf{\indent Returns:\ }
true or false 



 

 Let $S$ be a $p$\texttt{\symbol{45}}group, and let $L$ be a list of proto\texttt{\symbol{45}}essential subgroups of $S$. The operation \texttt{SemiCharacteristicGeneratedTest} checks that the subgroup 
\[\langle [O^p(\mathrm{Aut}(S)), S], Q_E \mid E \in L \rangle = S,\]
 where $Q_E$ is the semi\texttt{\symbol{45}}characteristic closure of $[N_S(E), E]$ in $E$. 

 This test applies at any prime $p$. It is based on \cite[Lemma 2.2(f)]{AOV17}. 

 }

 
\begin{Verbatim}[commandchars=!@|,fontsize=\small,frame=single,label=Example]
  !gapprompt@gap>| !gapinput@S := SmallGroup(2^7, 931);|
  <pc group of size 128 with 7 generators>
  !gapprompt@gap>| !gapinput@L := AllProtoEssentials(S);|
  [ Group([ f7, f1, f2, f4, f5, f6 ]), Group([ f5, f6, f7, f1, f3, f4 ]), 
  Group([ f5, f6, f7, f2, f3, f4 ]), Group([ f1, f2, f4, f6, f7 ]), 
  Group([ f1, f3, f5, f6, f7 ]) ]
  !gapprompt@gap>| !gapinput@SemiCharacteristicGeneratedTest(S, L{[1..3]});|
  true
  !gapprompt@gap>| !gapinput@SemiCharacteristicGeneratedTest(S, L{[3..5]});|
  false
\end{Verbatim}
 

\subsection{\textcolor{Chapter }{CoreTest (for IsPGroup, IsList)}}
\logpage{[ 4, 2, 12 ]}\nobreak
\hyperdef{L}{X8335DE5780AAD01F}{}
{\noindent\textcolor{FuncColor}{$\triangleright$\enspace\texttt{CoreTest({\mdseries\slshape S, L})\index{CoreTest@\texttt{CoreTest}!for IsPGroup, IsList}
\label{CoreTest:for IsPGroup, IsList}
}\hfill{\scriptsize (operation)}}\\
\textbf{\indent Returns:\ }
true or false 



 

 Let $S$ be a $p$\texttt{\symbol{45}}group, and let $L$ be a list of proto\texttt{\symbol{45}}essential subgroups of $S$. The operation \texttt{CoreTest} checks that if $N$ is a normal subgroup of $S$ such that $N$ is semi\texttt{\symbol{45}}characteristic (i.e. fixed by $O^p(O^{p'}(\mathrm{Aut}(E)))$) in each $E \in L$, then $N = 1$. 

 This test applies at any prime $p$. 

 }

 
\begin{Verbatim}[commandchars=!@|,fontsize=\small,frame=single,label=Example]
  !gapprompt@gap>| !gapinput@S := SmallGroup(256, 6366);|
  <pc group of size 256 with 8 generators>
  !gapprompt@gap>| !gapinput@L := AllProtoEssentials(S);|
  [ Group([ f1, f2, f4, f5, f6, f7, f8 ]), Group([ f7, f8, f3, f4, f5, f6 ]) ]
  !gapprompt@gap>| !gapinput@C := List(L, CharacteristicSubgroups);;|
  !gapprompt@gap>| !gapinput@Intersection(C);|
  [ Group([  ]) ]
  !gapprompt@gap>| !gapinput@CoreTest(S, L);|
  false
\end{Verbatim}
 }

 }

   
\chapter{\textcolor{Chapter }{Automizer Sequences}}\label{Chapter_AutomizerSequences}
\logpage{[ 5, 0, 0 ]}
\hyperdef{L}{X7E5B05DD7D971312}{}
{
  

 Let $S$ be a $p$\texttt{\symbol{45}}group, with proto\texttt{\symbol{45}}essential subgroups $E_1, \dots, E_n$. An \emph{automizer sequence} is a list $P := [H_{i_1}, \dots, H_{i_k}]$, where $H_{i_j} \leq \textrm{Aut}(E_{i_j})$ is a valid automizer for $E_{i_j}$ in $S$. In other words, $\textrm{Aut}_S(E_{i_j})$ is a Sylow $p$\texttt{\symbol{45}}subgroup of $H_{i_j}$, and $H_{i_j}/\textrm{Inn}(E_{i_j})$ contains a strongly $p$\texttt{\symbol{45}}embedded subgroup. Note that we do not include a choice of $H_S \leq \textrm{Aut}(S)$ here. These are prototypes for saturated fusion systems. In this chapter, we
describe processes to determine many important properties about automizer
sequences that will allow us to determine whether the corresponding fusion
system would be reduced. 
\section{\textcolor{Chapter }{The functions}}\label{Chapter_Automizer_Sequences_Section_The_functions}
\logpage{[ 5, 1, 0 ]}
\hyperdef{L}{X8173135A7D187358}{}
{
  

\subsection{\textcolor{Chapter }{FindReducedSequences (for IsPGroup, IsList)}}
\logpage{[ 5, 1, 1 ]}\nobreak
\hyperdef{L}{X8470CDE284146146}{}
{\noindent\textcolor{FuncColor}{$\triangleright$\enspace\texttt{FindReducedSequences({\mdseries\slshape S, L})\index{FindReducedSequences@\texttt{FindReducedSequences}!for IsPGroup, IsList}
\label{FindReducedSequences:for IsPGroup, IsList}
}\hfill{\scriptsize (operation)}}\\
\textbf{\indent Returns:\ }
a list of lists of integers 



 Let $S$ be a $p$\texttt{\symbol{45}}group, and let $L$ be a list of subgroups of $S$, all of which are proto\texttt{\symbol{45}}essential in $S$. We compute all possible automizer sequences involving the subgroups in $L$ that could lead to reduced fusion systems. Each sequence is returned as a list $A$ of integers, with one entry for each subgroup $E_i$ in $L$: $A[i] = 0$ if $E_i$ is not essential, and otherwise $A[i]$ is the position of $\textrm{Aut}_\mathcal{F}(E_i)$ in \texttt{ValidAutomizers(S, E{\textunderscore}i)}. }

 

\subsection{\textcolor{Chapter }{IsomorphismSequences (for IsPGroup, IsList, IsList)}}
\logpage{[ 5, 1, 2 ]}\nobreak
\hyperdef{L}{X833D9B1081E67538}{}
{\noindent\textcolor{FuncColor}{$\triangleright$\enspace\texttt{IsomorphismSequences({\mdseries\slshape S, P1, P2})\index{IsomorphismSequences@\texttt{IsomorphismSequences}!for IsPGroup, IsList, IsList}
\label{IsomorphismSequences:for IsPGroup, IsList, IsList}
}\hfill{\scriptsize (operation)}}\\
\textbf{\indent Returns:\ }
an automorphism of $S$, or \texttt{fail} 



 Given two automizer sequences $P_1$ and $P_2$, check whether there exists some $\phi \in \textrm{Aut}(S)$ such that $P_1^\phi = P_2$. This computation mimics isomorphism of fusion systems. The action of $\phi$ depends on the ordering of the $P_i$'s. }

 

\subsection{\textcolor{Chapter }{ValidAutomizers (for IsPGroup, IsPGroup)}}
\logpage{[ 5, 1, 3 ]}\nobreak
\hyperdef{L}{X7883D0967BF7D1EE}{}
{\noindent\textcolor{FuncColor}{$\triangleright$\enspace\texttt{ValidAutomizers({\mdseries\slshape S, E})\index{ValidAutomizers@\texttt{ValidAutomizers}!for IsPGroup, IsPGroup}
\label{ValidAutomizers:for IsPGroup, IsPGroup}
}\hfill{\scriptsize (operation)}}\\
\textbf{\indent Returns:\ }
a list 



 Finds valid automizers for $E$ in $S$. This is not part of the proto\texttt{\symbol{45}}essentials test and forms
the first step in the actual construction of fusion systems. Nonetheless, if a
subgroup is proto\texttt{\symbol{45}}essential, then it must have a valid
automizer. This function does not construct all subgroups of $\textrm{Aut}(E)/O_p(\textrm{Aut}(E))$, and instead uses iterative maximal subgroup condition to do so. }

 

\subsection{\textcolor{Chapter }{FindSmallestBorel (for IsPGroup, IsGroupOfAutomorphismsFiniteGroup)}}
\logpage{[ 5, 1, 4 ]}\nobreak
\hyperdef{L}{X7BE2D5B082F2BCA7}{}
{\noindent\textcolor{FuncColor}{$\triangleright$\enspace\texttt{FindSmallestBorel({\mdseries\slshape S, AutFE})\index{FindSmallestBorel@\texttt{FindSmallestBorel}!for IsPGroup, IsGroupOfAutomorphismsFiniteGroup}
\label{FindSmallestBorel:for IsPGroup, IsGroupOfAutomorphismsFiniteGroup}
}\hfill{\scriptsize (operation)}}\\
\textbf{\indent Returns:\ }
a group 



 Let $S$ be a $p$\texttt{\symbol{45}}group, and let $E < S$ with some fixed $H_E \leq \textrm{Aut}(E)$ part of an automizer sequence. The operation \texttt{FindSmallestBorel} computes the subgroup of $\textrm{Aut}(S)$ consisting of extensions from $N_{H_E}(\textrm{Aut}_S(E))$. If not every map in the normalizer subgroup exists, we return \texttt{fail}. }

 

\subsection{\textcolor{Chapter }{IsAutomizerBorelCompatible (for IsPGroup, IsGroupOfAutomorphismsFiniteGroup, IsGroupOfAutomorphismsFiniteGroup)}}
\logpage{[ 5, 1, 5 ]}\nobreak
\hyperdef{L}{X81A3ED2D7EBFAE5D}{}
{\noindent\textcolor{FuncColor}{$\triangleright$\enspace\texttt{IsAutomizerBorelCompatible({\mdseries\slshape S, AutFS, AutFE})\index{IsAutomizerBorelCompatible@\texttt{IsAutomizerBorelCompatible}!for IsPGroup, IsGroupOfAutomorphismsFiniteGroup, IsGroupOfAutomorphismsFiniteGroup}
\label{IsAutomizerBorelCompatible:for IsPGroup, IsGroupOfAutomorphismsFiniteGroup, IsGroupOfAutomorphismsFiniteGroup}
}\hfill{\scriptsize (operation)}}\\
\textbf{\indent Returns:\ }
true or false 



 Let $S$ be a $p$\texttt{\symbol{45}}group, and let $E < S$ with some fixed $H_E \leq \textrm{Aut}(E)$ part of an automizer sequence. Let $H_S \leq \textrm{Aut}(S)$ be a subgroup. The operation \texttt{IsAutomizerBorelCompatible} checks whether the restrictions of $H_S$ are \emph{compatible} with $H_E$. In particular, if $A_0 \leq \textrm{Aut}(E)$ consists of restrictions from $H_S$, we have that 
\[O^{p'}(\langle A_0, H_E \rangle) = H_E.\]
 }

 

\subsection{\textcolor{Chapter }{IsCompatibleSequence (for IsPGroup, IsList)}}
\logpage{[ 5, 1, 6 ]}\nobreak
\hyperdef{L}{X78334C0C8079F141}{}
{\noindent\textcolor{FuncColor}{$\triangleright$\enspace\texttt{IsCompatibleSequence({\mdseries\slshape S, P})\index{IsCompatibleSequence@\texttt{IsCompatibleSequence}!for IsPGroup, IsList}
\label{IsCompatibleSequence:for IsPGroup, IsList}
}\hfill{\scriptsize (operation)}}\\
\textbf{\indent Returns:\ }
true or false 



 Let $S$ be a $p$\texttt{\symbol{45}}group and let $P$ be an automizer sequence. The operation \texttt{IsCompatibleSequence} checks whether $P$ can lead to a saturated fusion system. Note that there can be false positives
\texttt{\symbol{45}}\texttt{\symbol{45}} $P$ may not lead to a saturated fusion system (based on the analysis of centric,
radical subgroups) but passes this test. }

 }

 }

   
\chapter{\textcolor{Chapter }{Other functions}}\label{Chapter_Other_functions}
\logpage{[ 6, 0, 0 ]}
\hyperdef{L}{X87C3D1B984960984}{}
{
  
\section{\textcolor{Chapter }{Group actions}}\label{Chapter_Other_functions_Section_Group_actions}
\logpage{[ 6, 1, 0 ]}
\hyperdef{L}{X87115591851FB7F4}{}
{
  

\subsection{\textcolor{Chapter }{OnQuotient}}
\logpage{[ 6, 1, 1 ]}\nobreak
\hyperdef{L}{X82E23FF785C4632A}{}
{\noindent\textcolor{FuncColor}{$\triangleright$\enspace\texttt{OnQuotient({\mdseries\slshape q})\index{OnQuotient@\texttt{OnQuotient}}
\label{OnQuotient}
}\hfill{\scriptsize (function)}}\\
\textbf{\indent Returns:\ }
a function 



 

 Let $G$ be a finite group, and $A$ be a normal subgroup of $G$. Let $q \colon G \to G/A$. The function \texttt{OnQuotient(q)} returns a function, such that for $t \in G/A$ and $x \in G$, \texttt{OnQuotient(q)(t, x)} returns $q(x^{-1}hx)$, where $h \in G$ satisfies $q(h) = t$. 

 }

 

\subsection{\textcolor{Chapter }{CentralizerMod (for IsGroup, IsGroup, IsGroup)}}
\logpage{[ 6, 1, 2 ]}\nobreak
\hyperdef{L}{X7FABB45881ECF282}{}
{\noindent\textcolor{FuncColor}{$\triangleright$\enspace\texttt{CentralizerMod({\mdseries\slshape G, A, B})\index{CentralizerMod@\texttt{CentralizerMod}!for IsGroup, IsGroup, IsGroup}
\label{CentralizerMod:for IsGroup, IsGroup, IsGroup}
}\hfill{\scriptsize (operation)}}\\
\textbf{\indent Returns:\ }
a subgroup of G 



 

 Let $G$ be a finite group, with subgroups $B \leq A \leq G$, where $B$ is normal in $G$. The operation \texttt{CentralizerMod(G, A, B)} returns the subgroup $C_G(A/B)$, given by: 

 
\[C_G(A/B) := \langle x \in G \mid [A, x] \leq B \rangle.\]
 

 }

 
\begin{Verbatim}[commandchars=!@|,fontsize=\small,frame=single,label=Example]
  !gapprompt@gap>| !gapinput@S := SmallGroup(5^6, 643);|
  <pc group of size 15625 with 6 generators>
  !gapprompt@gap>| !gapinput@U := UpperCentralSeries(S);|
  [ Group([ a6, a5, a4, a3, a1, a2 ]), Group([ a6, a5, a4, a3 ]), 
  Group([ a6, a5, a4 ]), Group([ a6, a5 ]), Group([ a6 ]), Group([  ]) ]
  !gapprompt@gap>| !gapinput@R := CentralizerMod(S, U[4], U[6]);|
  Group([ <identity> of ..., a2, a3, a4, a5, a6 ])
  !gapprompt@gap>| !gapinput@R = Centralizer(S, U[4]);|
  true
  !gapprompt@gap>| !gapinput@A := CentralizerMod(S, U[3], U[5]);|
  Group([ a1, <identity> of ..., a3, a4, a5, a6 ])
  !gapprompt@gap>| !gapinput@A = R;|
  false
\end{Verbatim}
 

\subsection{\textcolor{Chapter }{OnImage}}
\logpage{[ 6, 1, 3 ]}\nobreak
\hyperdef{L}{X7E7FE55986BF5133}{}
{\noindent\textcolor{FuncColor}{$\triangleright$\enspace\texttt{OnImage({\mdseries\slshape x, phi})\index{OnImage@\texttt{OnImage}}
\label{OnImage}
}\hfill{\scriptsize (function)}}\\
\textbf{\indent Returns:\ }
an element or a subset 



 

 Let $\phi$ be a map and let $x$ be either an element or a subset of the domain of $\phi$. The function \texttt{OnImage(x, phi)} returns $\phi(x)$. 

 }

 

\subsection{\textcolor{Chapter }{OnImageNM}}
\logpage{[ 6, 1, 4 ]}\nobreak
\hyperdef{L}{X7A23D935862598AE}{}
{\noindent\textcolor{FuncColor}{$\triangleright$\enspace\texttt{OnImageNM({\mdseries\slshape n})\index{OnImageNM@\texttt{OnImageNM}}
\label{OnImageNM}
}\hfill{\scriptsize (function)}}\\
\textbf{\indent Returns:\ }
a function 



 

 Let $G$ be a group and $n$ a nice monomorphism for (a subgroup of) $\mathrm{Aut}(G)$, with image $A$. The function \texttt{OnImageNM(n)} returns a function such that, for some element $\phi \in A$ and an element (or a subgroup) $x$ of $G$, \texttt{OnImageNM(n)(x, phi)} return $\hat{\phi}(x)$, where $\hat{\phi} \in \mathrm{Aut}(G)$ is the preimage of $\phi$ with respect to the injective map $n$. 

 }

 
\begin{Verbatim}[commandchars=!@|,fontsize=\small,frame=single,label=Example]
  !gapprompt@gap>| !gapinput@S := SmallGroup(5^6, 642);|
  <pc group of size 15625 with 6 generators>
  !gapprompt@gap>| !gapinput@AutS := AutomorphismGroup(S);|
  <group of size 1562500 with 10 generators>
  !gapprompt@gap>| !gapinput@n := NiceMonomorphism(AutS);|
  <action isomorphism>
  !gapprompt@gap>| !gapinput@AutS := n(AutS);|
  <permutation group of size 1562500 with 10 generators>
  !gapprompt@gap>| !gapinput@V := Group(S.2, S.6);|
  <pc group with 2 generators>
  !gapprompt@gap>| !gapinput@A := Stabilizer(AutS, V, OnImageNM(n));|
  <permutation group of size 12500 with 7 generators>
\end{Verbatim}
 

\subsection{\textcolor{Chapter }{OnImageTuples}}
\logpage{[ 6, 1, 5 ]}\nobreak
\hyperdef{L}{X87D74FC97C12B4D4}{}
{\noindent\textcolor{FuncColor}{$\triangleright$\enspace\texttt{OnImageTuples({\mdseries\slshape L, phi})\index{OnImageTuples@\texttt{OnImageTuples}}
\label{OnImageTuples}
}\hfill{\scriptsize (function)}}\\
\textbf{\indent Returns:\ }
a list 



 

 Let $\phi$ be a map and let $L$ be a list of elements in the domain of $\phi$. The function \texttt{OnImageTuples(L, phi)} returns a list containing $\phi(x)$ for each $x \in L$, while maintaining the order. 

 }

 

\subsection{\textcolor{Chapter }{OnImageTuplesNM}}
\logpage{[ 6, 1, 6 ]}\nobreak
\hyperdef{L}{X7B8996CD781FBFB9}{}
{\noindent\textcolor{FuncColor}{$\triangleright$\enspace\texttt{OnImageTuplesNM({\mdseries\slshape n})\index{OnImageTuplesNM@\texttt{OnImageTuplesNM}}
\label{OnImageTuplesNM}
}\hfill{\scriptsize (function)}}\\
\textbf{\indent Returns:\ }
a function 



 

 Let $G$ be a group and $n$ a nice monomorphism for (a subgroup of) $\mathrm{Aut}(G)$, with image $A$. The function \texttt{OnImageTuplesNM(n)} returns a function such that, for some element $\phi \in A$ and a list $L$ of elements in $G$, \texttt{OnImageTuplesNM(n)(L, phi)} return a list of elements $\hat{\phi}(x)$, where $x \in L$ and $\hat{\phi} \in \mathrm{Aut}(G)$ is the preimage of $\phi$ with respect to the injective map $n$. 

 }

 

\subsection{\textcolor{Chapter }{OnPreImage}}
\logpage{[ 6, 1, 7 ]}\nobreak
\hyperdef{L}{X82DDC150821611FA}{}
{\noindent\textcolor{FuncColor}{$\triangleright$\enspace\texttt{OnPreImage({\mdseries\slshape q})\index{OnPreImage@\texttt{OnPreImage}}
\label{OnPreImage}
}\hfill{\scriptsize (function)}}\\
\textbf{\indent Returns:\ }
a function 



 

 Let $A$ be a group that acts on the set $\Omega$ with kernel $K$, and let $q \colon A \to G$ be a group homomorphism with $\ker(q) \leq K$. The function \texttt{OnPreImage(q)} returns a function, such that, for some element $g \in G$ and $x \in \Omega$, \texttt{OnPreImage(q)(x, g)} returns $x^a$, where $a \in A$ is the preimage of $g$ with respect to $q$. 

 }

 

\subsection{\textcolor{Chapter }{OnPreImageTuples}}
\logpage{[ 6, 1, 8 ]}\nobreak
\hyperdef{L}{X7E32A67B83B9FBF6}{}
{\noindent\textcolor{FuncColor}{$\triangleright$\enspace\texttt{OnPreImageTuples({\mdseries\slshape q})\index{OnPreImageTuples@\texttt{OnPreImageTuples}}
\label{OnPreImageTuples}
}\hfill{\scriptsize (function)}}\\
\textbf{\indent Returns:\ }
a function 



 

 Let $A$ be a group that acts on the set $\Omega$ with kernel $K$, and let $q \colon A \to G$ be a group homomorphism with $\ker(q) \leq K$. The function \texttt{OnPreImageTuples(q)} returns a function, such that, for some element $g \in G$ and a subset $L \subseteq \Omega$, \texttt{OnPreImageTuples(q)(L, g)} returns a list consisting of $x^a$ for $x \in L$, where $a \in A$ is the preimage of $g$ with respect to $q$. 

 }

 }

 
\section{\textcolor{Chapter }{Group automorphisms}}\label{Chapter_Other_functions_Section_Group_automorphisms}
\logpage{[ 6, 2, 0 ]}
\hyperdef{L}{X783030917CB43959}{}
{
  

\subsection{\textcolor{Chapter }{Automizer (for IsGroup, IsGroup)}}
\logpage{[ 6, 2, 1 ]}\nobreak
\label{AutoDoc_generated_group3}
\hyperdef{L}{X7D9EC37E869F3256}{}
{\noindent\textcolor{FuncColor}{$\triangleright$\enspace\texttt{Automizer({\mdseries\slshape G, H})\index{Automizer@\texttt{Automizer}!for IsGroup, IsGroup}
\label{Automizer:for IsGroup, IsGroup}
}\hfill{\scriptsize (operation)}}\\
\noindent\textcolor{FuncColor}{$\triangleright$\enspace\texttt{Automizer({\mdseries\slshape A, H})\index{Automizer@\texttt{Automizer}!for IsGroupOfAutomorphismsFiniteGroup, IsGroup and IsFinite}
\label{Automizer:for IsGroupOfAutomorphismsFiniteGroup, IsGroup and IsFinite}
}\hfill{\scriptsize (operation)}}\\
\textbf{\indent Returns:\ }
a subgroup of Aut(H) 



 

 Let $G$ and $H$ be subgroups of some finite group. The group \texttt{Automizer(G, H)} is the automizer subgroup $\mathrm{Aut}_G(H)$, given by: 

 
\[\mathrm{Aut}_G(H) := \langle c_g \mid g \in N_G(H) \rangle,\]
 

 where $c_g \colon H \to H$ is the conjugation automorphism. It is a subgroup of $\mathrm{Aut}(H)$. 

 }

 
\begin{Verbatim}[commandchars=!@|,fontsize=\small,frame=single,label=Example]
  !gapprompt@gap>| !gapinput@G := AlternatingGroup(6);|
  Alt( [ 1 .. 6 ] )
  !gapprompt@gap>| !gapinput@S := SylowSubgroup(G, 2);|
  Group([ (1,2)(3,4), (1,3)(2,4), (1,2)(5,6) ])
  !gapprompt@gap>| !gapinput@M := MaximalSubgroups(S);|
  [ Group([ (3,4)(5,6), (1,2)(3,4) ]), Group([ (1,3)(2,4), (1,2)(3,4) ]), 
  Group([ (1,3,2,4)(5,6), (1,2)(3,4) ]) ]
  !gapprompt@gap>| !gapinput@AutGM := List(M, A -> Automizer(G, A));|
  [ <group with 2 generators>, <group with 2 generators>, 
  <group with 2 generators> ]
  !gapprompt@gap>| !gapinput@List(AutGM, StructureDescription);|
  [ "S3", "S3", "C2" ]
\end{Verbatim}
 Alternatively, if $A$ is a subgroup of $\mathrm{Aut}(G)$, then \texttt{Automizer(A, H)} returns the subgroup $\mathrm{Aut}_A(H)$ of $\mathrm{Aut}(H)$ given by: 

 
\[\mathrm{Aut}_A(H) := \langle \alpha|_H \mid \alpha \in A, \alpha(H) = H
\rangle.\]
 

 
\begin{Verbatim}[commandchars=!@|,fontsize=\small,frame=single,label=Example]
  !gapprompt@gap>| !gapinput@S := SmallGroup(5^6, 642);|
  <pc group of size 15625 with 6 generators>
  !gapprompt@gap>| !gapinput@V := Group(S.2, S.6);|
  <pc group with 2 generators>
  !gapprompt@gap>| !gapinput@AutS := AutomorphismGroup(S);|
  <group of size 1562500 with 10 generators>
  !gapprompt@gap>| !gapinput@AutSV := Automizer(AutS, V);|
  <group with 7 generators>
  !gapprompt@gap>| !gapinput@AutV := AutomorphismGroup(V);;|
  !gapprompt@gap>| !gapinput@n := NiceMonomorphism(AutV);|
  <action isomorphism>
  !gapprompt@gap>| !gapinput@AutV := n(AutV);;|
  !gapprompt@gap>| !gapinput@Size(AutV);|
  480
  !gapprompt@gap>| !gapinput@AutSV := n(AutSV);|
  <permutation group with 7 generators>
  !gapprompt@gap>| !gapinput@Index(AutV, AutSV);|
  24
  !gapprompt@gap>| !gapinput@Size(AutSV);|
  20
\end{Verbatim}
 
\subsection{\textcolor{Chapter }{Restriction of homomorphisms}}\label{AutoDoc_generated_group4}
\logpage{[ 6, 2, 2 ]}
\hyperdef{L}{X874E7F578433266D}{}
{
\noindent\textcolor{FuncColor}{$\triangleright$\enspace\texttt{RestrictedHomomorphismOnto({\mdseries\slshape phi, X})\index{RestrictedHomomorphismOnto@\texttt{RestrictedHomomorphismOnto}!for IsGroupHomomorphism, IsGroup}
\label{RestrictedHomomorphismOnto:for IsGroupHomomorphism, IsGroup}
}\hfill{\scriptsize (operation)}}\\
\noindent\textcolor{FuncColor}{$\triangleright$\enspace\texttt{RestrictedAutomorphism({\mdseries\slshape phi, X})\index{RestrictedAutomorphism@\texttt{RestrictedAutomorphism}!for IsGroupHomomorphism and IsBijective, IsGroup}
\label{RestrictedAutomorphism:for IsGroupHomomorphism and IsBijective, IsGroup}
}\hfill{\scriptsize (operation)}}\\
\textbf{\indent Returns:\ }
a group automorphism 



 

 Let $\phi \colon A \to B$ be a group homomorphism, and let $X \leq A$. The operation \texttt{RestrictedHomomorphismOnto(phi, X)} returns the surjective group homomorphism $\phi|_X \colon X \to \phi(X)$ given by restriction. If $\phi$ is also an isomorphism that fixes $X$, then the operation \texttt{RestrictedAutomorphism(phi, X)} returns the group automorphism $\phi|_X \colon X \to \phi(X)$ given by restriction. 

 }

 
\subsection{\textcolor{Chapter }{Transfer homomorphism}}\label{AutoDoc_generated_group5}
\logpage{[ 6, 2, 3 ]}
\hyperdef{L}{X80F6FD3E7C7E4E8D}{}
{
\noindent\textcolor{FuncColor}{$\triangleright$\enspace\texttt{TransferHomomorphism({\mdseries\slshape G, S})\index{TransferHomomorphism@\texttt{TransferHomomorphism}!for IsGroup, IsGroup}
\label{TransferHomomorphism:for IsGroup, IsGroup}
}\hfill{\scriptsize (operation)}}\\
\noindent\textcolor{FuncColor}{$\triangleright$\enspace\texttt{TransferHomomorphismFp({\mdseries\slshape G, S})\index{TransferHomomorphismFp@\texttt{TransferHomomorphismFp}!for IsGroup, IsPGroup}
\label{TransferHomomorphismFp:for IsGroup, IsPGroup}
}\hfill{\scriptsize (operation)}}\\
\textbf{\indent Returns:\ }
a homomorphism 



 

 Let $G$ be a finite group and $S \leq G$. The operation \texttt{TransferHomomorphism(G, S)} returns the transfer homomorphism $\mathrm{tr}_S^G \colon G \to S/S'$. Moreover, if $S$ is a $p$\texttt{\symbol{45}}group, then the operation \texttt{TransferHomomorphismFp(G, S)} returns the transfer homomorphism $\mathrm{tr'}_S^G \colon G \to S/\Phi(S)$. 

 }

 }

 
\section{\textcolor{Chapter }{Semi\texttt{\symbol{45}}characteristic subgroups}}\label{Chapter_Other_functions_Section_Semi-characteristic_subgroups}
\logpage{[ 6, 3, 0 ]}
\hyperdef{L}{X7EA0263C874496CA}{}
{
  Let $S$ be a $p$\texttt{\symbol{45}}group and $A \leq S$. We say that $A$ is \emph{semi\texttt{\symbol{45}}characteristic} in $S$ if 
\begin{itemize}
\item  $A$ is normal in $S$, and 
\item  $A$ is fixed by $O^p(O^{p'}(\mathrm{Aut}(S)))$. 
\end{itemize}
 

\subsection{\textcolor{Chapter }{SemiAutomorphismGroup (for IsPGroup)}}
\logpage{[ 6, 3, 1 ]}\nobreak
\hyperdef{L}{X7E8444CF8156B41C}{}
{\noindent\textcolor{FuncColor}{$\triangleright$\enspace\texttt{SemiAutomorphismGroup({\mdseries\slshape S})\index{SemiAutomorphismGroup@\texttt{SemiAutomorphismGroup}!for IsPGroup}
\label{SemiAutomorphismGroup:for IsPGroup}
}\hfill{\scriptsize (attribute)}}\\
\textbf{\indent Returns:\ }
a subgroup of Aut(S) 



 

 Let $S$ be a $p$\texttt{\symbol{45}}group. The operation \texttt{SemiAutomorphismGroup} returns the semi\texttt{\symbol{45}}automorphism group of $S$, which is given by $\langle \mathrm{Inn}(S), O^p(O^{p'}(\mathrm{Aut}(S)))\rangle$. 

 }

 

\subsection{\textcolor{Chapter }{SemiCharacteristicSubgroups (for IsPGroup)}}
\logpage{[ 6, 3, 2 ]}\nobreak
\hyperdef{L}{X7E88D38982158EB1}{}
{\noindent\textcolor{FuncColor}{$\triangleright$\enspace\texttt{SemiCharacteristicSubgroups({\mdseries\slshape S})\index{SemiCharacteristicSubgroups@\texttt{SemiCharacteristicSubgroups}!for IsPGroup}
\label{SemiCharacteristicSubgroups:for IsPGroup}
}\hfill{\scriptsize (attribute)}}\\
\textbf{\indent Returns:\ }
a list 



 

 Let $S$ be a $p$\texttt{\symbol{45}}group. The attribute \texttt{SemiCharacteristicSubgroups} returns all the semi\texttt{\symbol{45}}characteristic subgroups of $S$. These are the normal subgroups of $S$ fixed by $O^p(O^{p'}(\mathrm{Aut}(S)))$. 

 }

 

\subsection{\textcolor{Chapter }{IsSemiCharacteristic (for IsPGroup, IsPGroup)}}
\logpage{[ 6, 3, 3 ]}\nobreak
\hyperdef{L}{X840438E7843A6128}{}
{\noindent\textcolor{FuncColor}{$\triangleright$\enspace\texttt{IsSemiCharacteristic({\mdseries\slshape S, A})\index{IsSemiCharacteristic@\texttt{IsSemiCharacteristic}!for IsPGroup, IsPGroup}
\label{IsSemiCharacteristic:for IsPGroup, IsPGroup}
}\hfill{\scriptsize (operation)}}\\
\textbf{\indent Returns:\ }
true or false 



 

 Let $S$ be a $p$\texttt{\symbol{45}}group and $A$ be a subgroup of $S$. The operation \texttt{IsSemiCharacteristic} checks whether $A$ is semi\texttt{\symbol{45}}characteristic in $S$. That is, we check whether: 

 
\begin{itemize}
\item  $A$ is normal in $S$, and 
\item  $A$ is fixed by $O^p(O^{p'}(\mathrm{Aut}(S)))$. 
\end{itemize}
 

 }

 

\subsection{\textcolor{Chapter }{SemiCharacteristicClosure (for IsPGroup, IsPGroup)}}
\logpage{[ 6, 3, 4 ]}\nobreak
\hyperdef{L}{X8759D1937F245ED7}{}
{\noindent\textcolor{FuncColor}{$\triangleright$\enspace\texttt{SemiCharacteristicClosure({\mdseries\slshape S, A})\index{SemiCharacteristicClosure@\texttt{SemiCharacteristicClosure}!for IsPGroup, IsPGroup}
\label{SemiCharacteristicClosure:for IsPGroup, IsPGroup}
}\hfill{\scriptsize (operation)}}\\
\textbf{\indent Returns:\ }
a subgroup of S 



 

 Let $S$ be a $p$\texttt{\symbol{45}}group and $A$ be a subgroup of $S$. The operation \texttt{SemiCharacteristicClosure} returns the smallest semi\texttt{\symbol{45}}characteristic subgroup of $S$ containing $A$. This is the intersection of all semi\texttt{\symbol{45}}characteristic
subgroups of $S$ containing $A$. 

 }

 
\begin{Verbatim}[commandchars=!@|,fontsize=\small,frame=single,label=Example]
  !gapprompt@gap>| !gapinput@S := SmallGroup(2^8, 6331);|
  <pc group of size 256 with 8 generators>
  !gapprompt@gap>| !gapinput@E1 := Group(S.1, S.2, S.4, S.5, S.6, S.7, S.8);|
  <pc group with 7 generators>
  !gapprompt@gap>| !gapinput@E2 := Group(S.1, S.2, S.4, S.5, S.7, S.8);|
  <pc group with 6 generators>
  !gapprompt@gap>| !gapinput@A := Group(S.5*S.8);|
  <pc group with 1 generator>
  !gapprompt@gap>| !gapinput@AutomorphismGroup(E1);|
  <group of size 6144 with 12 generators>
  !gapprompt@gap>| !gapinput@AutomorphismGroup(E2);|
  <group of size 12288 with 13 generators>
  !gapprompt@gap>| !gapinput@IsCharacteristicSubgroup(E1, A);|
  false
  !gapprompt@gap>| !gapinput@IsCharacteristicSubgroup(E2, A);|
  false
  !gapprompt@gap>| !gapinput@IsPGroup(E1);|
  true
  !gapprompt@gap>| !gapinput@IsPGroup(E2);|
  true
  !gapprompt@gap>| !gapinput@SemiAutomorphismGroup(E1);|
  <group of size 384 with 8 generators>
  !gapprompt@gap>| !gapinput@SemiAutomorphismGroup(E2);|
  <group of size 3072 with 11 generators>
  !gapprompt@gap>| !gapinput@IsPGroup(A);|
  true
  !gapprompt@gap>| !gapinput@IsSemiCharacteristic(E1, A);|
  true
  !gapprompt@gap>| !gapinput@IsSemiCharacteristic(E2, A);|
  false
  !gapprompt@gap>| !gapinput@SemiCharacteristicClosure(E2, A);|
  Group([ f4, f5, f8, f4*f5, f4*f8, f5*f8, f4*f5*f8 ])
\end{Verbatim}
 }

 }

 \def\bibname{References\logpage{[ "Bib", 0, 0 ]}
\hyperdef{L}{X7A6F98FD85F02BFE}{}
}

\bibliographystyle{alpha}
\bibliography{FusionSystems.bib}

\addcontentsline{toc}{chapter}{References}

\def\indexname{Index\logpage{[ "Ind", 0, 0 ]}
\hyperdef{L}{X83A0356F839C696F}{}
}

\cleardoublepage
\phantomsection
\addcontentsline{toc}{chapter}{Index}


\printindex

\immediate\write\pagenrlog{["Ind", 0, 0], \arabic{page},}
\newpage
\immediate\write\pagenrlog{["End"], \arabic{page}];}
\immediate\closeout\pagenrlog
\end{document}
